\documentclass[a4paper,12pt]{book}
\usepackage[a4paper, portrait, margin=0.5in]{geometry}
\usepackage[utf8]{inputenc}
\usepackage{graphicx}
\usepackage{braket}
\usepackage{dsfont}
\usepackage{amsmath}
\usepackage{mathtools}
\usepackage{cancel}
\usepackage{bbold}
\usepackage{titlepic}
\usepackage{amssymb}
\usepackage{amsthm}
\usepackage{soul}
\usepackage{hyperref}
\setcounter{secnumdepth}{3}
\setcounter{tocdepth}{3}
\setcounter{MaxMatrixCols}{20}

\newcommand{\limit}[2]{\underset{#1 \rightarrow #2}{\lim} \,}
\newcommand{\arrowlim}[2]{\overset{#1 \rightarrow #2}{\longrightarrow}}
\newcommand{\arrowlimtwo}[4]{\overset{\scriptsize{\begin{array}{cc}
				#1 {}&\rightarrow #2\\
				#3 &\rightarrow #4
		\end{array}}
	}{\longrightarrow}}

\newcommand{\Lagr}{\mathcal{L}}
\newcommand{\Four}{\mathcal{F}}
\newcommand{\TdZ}{\mathcal{Z}}
\newcommand{\R}{\mathbb{R}}
\newcommand{\C}{\mathbb{C}}
\newcommand{\N}{\mathbb{N}}
\newcommand{\Z}{\mathbb{Z}}
\newcommand{\nextpassage}{\nonumber \\ &\downarrow \nonumber \\}
\newcommand{\firstpassage}{ \\ &\downarrow \nonumber \\}
\newcommand{\spacer}{\quad ; \quad}
\newcommand{\except}[1]{\backslash \{ #1 \} }
\newcommand{\taleche}{\; : \;}
\newcommand{\bolditem}[1]{\item \textbf{#1}}
\newcommand{\separator}{\nonumber \\ 
	\nonumber \\}
\newcommand{\double}[2]{\left\{ \begin{array}{cc}
		#1\\
		#2
	\end{array}
	\right.
}
\newcommand{\triple}[3]{\left\{ \begin{array}{ccc}
		#1\\
		#2\\
		#3
	\end{array}
	\right.
}	
\newcommand{\quadruple}[4]{\left\{ \begin{array}{ccc}
		#1\\
		#2\\
		#3\\
		#4
	\end{array}
	\right.
}	

\newcommand{\nosgn}{\;\;\,}

\newcommand{\quadvec}[4]{\left( \begin{array}{ccc}
		#1\\
		#2\\
		#3\\
		#4
	\end{array}
	\right)
}

\newcommand{\dubvec}[2]{\left( \begin{array}{ccc}
		#1\\
		#2
	\end{array}
	\right)
}


\begin{document}

\author{Alessandro Marcelli}
\title{\centering
	Raccolta di esercizi e spunti di teoria per il corso di\newline
	Metodi e Modelli Matematici per la Fisica 2,\newline
	svolti da Davide Bufalini, Paolo Proia e Adriano Chialastri}
\date{2023}



\begin{titlepage}
	\vspace*{\fill}
	\begin{center}
		\scalebox{1.75}{Esercizi svolti da gente brava e spunti di teoria}\\
		\scalebox{1.75}{per la prova scritta del corso di}\\
		\scalebox{1.75}{Metodi e Modelli Matematici per la Fisica 2}\\[0.5cm]
		\begin{figure}[!htb]
			\centering
			\includegraphics[width=.5\textwidth]{./turtwig.jpg}
		\end{figure}
		\scalebox{1.25}{a cura di}\\
		\scalebox{1.25}{ }\\
		\scalebox{1.5}{Alessandro Marcelli}\\
		\scalebox{1.25}{ }\\
		\scalebox{1.25}{2023}
	\end{center}
	\vspace{\fill}	
\end{titlepage}

\clearpage
\thispagestyle{empty}

\frontmatter
\tableofcontents
\chapter{Note introduttive}

\section{Ciancio alle bande}

Questa ""dispensa"" è una raccolta di materiale utile (almeno al momento della scrittura di questa introduzione) alla prova scritta dell'esame  di Metodi e Modelli Matematici per la Fisica 2 del corso di laurea magistrale in Fisica dell'Univeristà Tor Vergata. 

Sono esercizi svolti da gente brava che, come esercizio personale ai fini di comprenderli al meglio, ho raccolto insieme, modificandoli quel tanto che bastava per passare dal "calcolo su foglio" a roba che sembra effettivamente una dispensa. Sul successo dell'operazione lascio il giudizio ai posteri, ammesso che ce ne saranno.

Quando possibile ho cercato di aggiungere dei minimi richiami di teoria qualora fossero assenti dagli appunti che nel corso del tempo ho \st{sciacallato} raccolto durante la preparazione di questa prova scritta.

Ben lungi dall'essere un lavoro perfetto, ma sinceramente spero di non fare in tempo a renderlo tale (RIP Equazioni Differenziali in campo complesso). 

\section{Crediti}

Nei titoli degli argomenti sono creditati i vari autori dei materiali raccolti. Qualora una sezione fosse composta da materiale di più persone, queste sono prima creditate nel titolo della sezione e poi nelle sottosezioni rilevanti. Se invece una sezione è farina del sacco di un solo autore, questo viene creditato solo nel titolo della sezione.

Detto in parole povere, ste dispense sono l'equivalente di un mixtape, e a parte che con il buon Paolo, non ho collaborato con gli altri due autori (dei quali ho però la benedizione a pubblicare queste dispense), mi sono limitato a fare l'amanuense con il loro lavoro. Ergo se avete cacciato soldi per avere sto documento vi hanno truffato.

\section{Contatti}

Qualora eventuali futuri \st{masochisti} utenti volessero contribuire a quest'accrocco mi si può raggiunere all'indirizzo \textbf{marcelli.alessandro.92@gmail.com}




\mainmatter
\chapter{Funzioni analitiche, prolungamenti ed espansioni asintotiche}

\section{Nozioni generali di Analisi Complessa (Alessandro Marcelli)}

\subsection{I numeri complessi}

I numeri complessi possono essere espressi come somma di una \textbf{pare reale} e di una \textbf{parte immaginaria}
\begin{align}
	z = x + iy \spacer x,y \in \R
\end{align}
Si definisce il \textbf{complesso coniugato} di un numero complesso come
\begin{align}
	\overline{z} = x-iy
\end{align}
Possiamo definire il suo \textbf{modulo} come
\begin{align}
	|z|^2 = z \cdot \overline{z} = x^2 + y^2
\end{align}
Utilizzando queste due definizioni si possono riscrivere le parti reale e immaginaria come
\begin{align}
	x = \frac{z+ \overline{z}}{2}\\
	y = \frac{z- \overline{z}}{2i}
\end{align}
Un numero complesso può essere rappresentato anche in \textbf{forma polare} nel seguente modo
\begin{align}
	&\double{x = \rho \cos (\theta)}{y = \rho \sin (\theta)} \spacer \double{\rho = \sqrt{x^2 + y^2}}{\theta = \arctan (\frac{y}{x})}\firstpassage
	&\double{z =  \rho (\cos (\theta) + i\sin(\theta))}{\overline{z} = \rho (\cos (\theta) - i\sin(\theta))}
\end{align}
Possiamo inoltre usare la \textbf{formula di Eulero} per riscrivere
\begin{align}
	&e^{i\theta} = \cos(\theta)+ i \sin (\theta)\firstpassage
	&\double{z = \rho e^{+i\theta}}{\overline{z} = \rho e^{-i\theta}}\\
	&\cos (\theta) = \frac{e^{i\theta} + e^{-i\theta}}{2} \spacer x= \frac{e^{i\theta} + e^{-i\theta}}{2\rho}\\
	&\sin (\theta) = \frac{e^{i\theta} - e^{-i\theta}}{2i} \spacer x= \frac{e^{i\theta} - e^{-i\theta}}{2i\rho}
\end{align}

\newpage

\subsection{Funzioni di variabile complessa}

Una funzione di variabile complessa $f(z)$ per definizione può essere scritta come
\begin{align}
	f(z) = u(x,y) + i v(x,y) \spacer x,y, u(x,y), v(x,y) \in \R
\end{align}
In analogo col caso reale, $f(z)$ si dice \textbf{differenziabile}\cite{MfP} in $z_0 \in \C$ se il limite
\begin{align}
	\limit{\Delta z}{0} \frac{f(z_0 + \Delta z) - f(z_0)}{\Delta z}
\end{align}
\begin{enumerate}
	\item esiste
	\item è finito
	\item il suo valore non cambia in funzione della direzione dal quale lo si approccia
\end{enumerate}
La difficoltà aggiuntiva rispetto al caso reale è dovuto al fatto che in campo complesso il numero di direzioni possibili è infinito. Le condizioni sono quindi molto più restrittive del caso reale, ovvero
\begin{enumerate}
	\item Devono essere rispettate le \textbf{condizioni di Cauchy-Riemann}:
	\begin{align}
		\frac{\partial u(x,y)}{\partial x} &= +\frac{\partial v(x,y)}{\partial y}\\
		\frac{\partial u(x,y)}{\partial y} &= -\frac{\partial v(x,y)}{\partial x}		
	\end{align}
	\item Le derivate parziali prime di $u(x,y)$ e $v(x,y)$ devono essere continue
\end{enumerate}

Le funzioni di variabile complessa si dividono in due categorie:
\begin{enumerate}
	\item \textbf{funzioni a singolo valore}, ovvero che assumono lo stesso valore in $z_0$ indipendentemente dal percorso fatto per raggiungerlo. Ad esempio
	\begin{align}
		f(z) = \frac{1}{z}
	\end{align}
	\item \textbf{funzioni multivalore}, ovvero il cui valore varia a seconda del percorso fatto per raggiungere $z_0$. Ad esempio
	\begin{align}
		g(z) = \ln(z) = \ln |z| + i \arg (z)
	\end{align}
\end{enumerate}

Se la funzione è differenziabile ed è a \textbf{singolo valore} per ogni punto di un insieme $\mathcal{D}$, allora si dice che è \textbf{analitica} in questo insieme, che prende il nome di \textbf{dominio di analiticità}. All'interso di $\mathcal{D}$ essa non può assumere massimi locali.

Una funzione analitica può essere anche riscritta come (RITROVA GLI APPUNTI DI SANTINI)
\begin{align}
	f(z) = u(x,y) + iv(x,y) \to f(z) = u(z,\overline{z}) + i v(z,\overline{z})
\end{align}

Se una funzione è analitica in un punto, questo si dice \textbf{punto regolare} della funzione, altrimenti prende il nome di \textbf{punto singolare}.

Per le funzioni multivalore spesso si può "barare" utilizzando le \textbf{superfici di Riemann}, che ci permettono di usare le comode proprietà delle funzioni analitiche.

Il \textbf{teorema di Cauchy} ci dice che, per curve chiuse $\gamma$
\begin{align}
	\oint_\gamma dz \; f(z) = 0 \quad \forall \gamma \in \mathcal{D}
\end{align}
Siccome qualunque curva chiusa può essere scritta come combinazioni di curve con gli stessi estremi, questo implica che fissati gli estremi di integrazione il valore dell'integrale è indipendente dal percorso.


\newpage
\subsection{Punti di non analiticità}

\subsubsection{Poli e zeri}

Trovare gli appunti di Santini

\subsubsection{Residui}

Se in un insieme $A$ la nostra $f(z)$ presenta singolarità isolate $z_i$, non vale più il teorema di Cauchy, e quindi non si annulla l'integrale, il cui valore viene dato dalla somma dei \textbf{residui della funzione} in tali punti
\begin{align}
	&\oint_\gamma dz \; f(z) = \sum_{i=1}^N \text{Res} (f(z), z_i)\\
	&\text{Res} (f(z), z_i) = \limit{z}{z_i} \frac{1}{(n-1)!} \left( \frac{d^{n-1}}{dz^{n-1}} \left[ (z-z_0)^n f(z) \right] \right) \spacer n = \text{ordine del polo}
\end{align}


\newpage

\section{Indicatore Logaritmico (Davide Bufalini)}
Data una funzione $f(z)$ meromorfa ,si definisce l'\textbf{indicatore logaritmico} come
\begin{align}
	L_f(z) = \frac{d}{dz}ln(f(z)) = \frac{f'(z)}{f(z)}
\end{align}
$L_f(z)$ sarà dotata di poli semplici.


Facciamo alcune osservazioni:
\begin{enumerate}
	\item Se $f(z)$ ha $N$ poli e $K$ zeri allora
	\begin{align}
		\oint_C \frac{dz}{2\pi i} L_f(z) = \sum_{i=1}^{N} m_i - \sum_{l=1}^{M}r_l
	\end{align}
	Le quantità $m_i$ e $r_l$ sono rispettivamente le molteplicità dei poli e degli zeri di $f(z)$.
	\item Per l'indicatore logaritmico vale la seguente relazione
	\begin{align}
		\int_{z_i}^{z_f}\frac{dz}{2\pi i} L_f(z) = \frac{1}{2\pi} \Delta_c Arg(f(z))
	\end{align}
	\item In generale vale che
	\begin{align}
		\oint_C \frac{dz}{2\pi i} L_f(z) \phi(z) = \sum_{i=1}^{N} m_i\phi(a_i) - \sum_{l=1}^{M}r_l\phi(b_lcc)
	\end{align}
	Dove $a_i$ e $b_l$ sono rispettivamente i poli e gli zeri di $\phi(z)$. 
\end{enumerate}

\newpage

\section{Inversione e reciprocità locale (Davide Bufalini)}

Sia una $f(z)$ analitica con centro dello sviluppo in $z_0$ tale che
\begin{align}
	f(z) = \sum_{n=0}^{+\infty} a_n (z-z_0)^n	\label{eq1}
\end{align}
allora, preso un intorno di $z_0$, vi si possino definire la reciproca della funzione $\frac{1}{f}$ e la sua inversa $z^{-1}$.

\subsection{Funzione reciproca}

Definiamo la reciproca come
\begin{align}
	\frac{1}{f} = \sum_{k} c_k (z-z_0)^k
\end{align}
Allora per definizione avremo, considerando la \ref{eq1}
\begin{align}
	&\sum_{k} c_k (z-z_0)^k \cdot \sum_{n=0} a_n (z-z_0)^n = 1\firstpassage
	&\sum_{k,n} c_ka_n (z-z_0)^{n+k} = 1 \label{eq2}
\end{align}
Riscriviamo la \ref{eq2} come
\begin{align}
	\sum_{l=0}^{+\infty} d_l (z-z_0)^l = 1
\end{align}
definendo
\begin{align}
	l &= n+k\\
	d_l &= \sum_{n=0}^{l}c_{l-n}a_n
\end{align}
Vediamo come
\begin{align}
	1 = d_0 &= c_0a_0 \quad \rightarrow& \quad &c_0 = \frac{1}{a_0}\\
	0 = d_1 &= c_0a_1 + c_1a_0 \quad \rightarrow& \quad &c_1 = -\frac{a_1}{a_0^2} = -\frac{a_1}{a_0} c_0\\
	0 = d_2 &= c_0a_2 + c_1a_1 + c_2a_0 \quad \rightarrow& \quad &c_2 = \dots = -\frac{1}{a_0}\left( -\frac{a_2}{a_0} + \frac{a_1^2}{a_0^2} \right) \\
	\dots \nonumber
\end{align}
Andiamo ora a scrivere
\begin{align}
	\frac{1}{f} &= \frac{1}{a_0}\cdot \frac{1}{1 + \frac{a_1}{a_0}(z-z_0) + \frac{a_2}{a_0} (z-z_0)^2 + \dots} := \frac{1}{a_0} \cdot \frac{1}{1-h(z)}
\end{align}
Siccome $|h(z)| \arrowlim{z}{z_0} 0$, possiamo appoggiarci al concetto di serie geometrica
\begin{align}
	\frac{1}{f} = \frac{1}{a_0}\cdot (1 + h(z) + h^2(z) + \dots)
\end{align}

\subsection{Funzione inversa}

Iniziamo notando come $F(z_0) = a_0$, da cui troviamo che, definendo
\begin{align}
	g(w) = f^{-1}(w) = \sum_{n=0}^{+\infty} b_n(w-w_0)^n
\end{align}
Si ottiene
\begin{align}
	g(w_0) = f^{-1}(w_0) = b_0 = z_0
\end{align}

Possiamo quindi procedere in due modi:
\begin{enumerate}
	\item \textbf{Per serie:}
	\begin{align}
		&f(z) - f(z_0) = \sum_{n=1}^{+\infty} a_n (z-z_0)^n= w-w_0\firstpassage
		& z -z_0 = g(w) - g(w_0) = 	 \sum_{n=1}^{+\infty} b_n(w-w_0)^n	\firstpassage
		& w - w_0 = \sum_{n=1}^{+\infty} a_n (\sum_{k=1}^{+\infty} b_k (w-w_0)^k)^n
	\end{align}
	\item \textbf{Formula di Lagrange:}
	\begin{align}
		b_0 &= g(w_0) = z_0\\
		b_n &= \frac{1}{n!} \frac{d^{n-1}}{dz^{n-1}} \left. \left[ \frac{z-z_0}{f(z) - f(z_0)} \right] \right|_{z=z_0}
	\end{align}
	
\end{enumerate}

\newpage

\subsection{Esempio: Numeri di Eulero}

Data la relazione
\begin{align}
	\frac{1}{\coth (z)} &= \sum_{n=0}^{+\infty} \frac{E_n}{n!}z^n\\
	E_{2n+1} &= 0 \quad \forall n\\
	E_n &= \text{numeri di eulero}
\end{align}

\subsection{Esempio: Numeri di Bernoulli}
\newpage


\section{Espansione di Mittag-Leffler (Davide Bufalini)}

Presa una $f(z)$ analitica nell'origine e con infiniti poli semplici che si accumulano all'infinito, vale la seguente approssimazione
\begin{align}
	f(z) &= f(0) + \sum_{n=1}^{+\infty} R_n \cdot \left( \frac{1}{z-z_n} + \frac{1}{z_n} \right)\\
	R_n &= Res(f, z_n)
\end{align}

\subsection{Esempio: cotangente}

\subsubsection{Calcolo dell'espansione}

La funzione  
\begin{align}
	f(z) = \cot(z)
\end{align}
ha poli in $z_n = n\pi$ con $n\in \Z$, ma \textbf{non} è analitica in $z=0$.

Per eliminare la singolarità in 0 considero allora la funzione
\begin{align}
	f(z) = \cot(z) - \frac{1}{z}
\end{align}
I cui residui saranno
\begin{align}
	R_n = \limit{z}{n\pi} \left[ (z-n\pi) \cdot \frac{z\cot(z) -1}{z} \right] = \dots = 1  
\end{align}
Applicando lo sviluppo ottengo
\begin{align}
	\cot(z) - \frac{1}{z} &= \sum_{n\neq 0} \left( \frac{1}{z-n\pi} + \frac{1}{n\pi}\right) \firstpassage
	\cot(z) &= \frac{1}{z} + \sum_{n\neq 0} \left( \frac{1}{z-n\pi} + \frac{1}{n\pi}\right) = \nonumber\\
	&= \frac{1}{z} + \sum_{n>0} \left[ \frac{1}{z-n\pi} + \frac{1}{n\pi} - \left( \frac{1}{-z-n\pi} + \frac{1}{n\pi} \right) \right] = \nonumber\\
	&= \dots  = \nonumber\\
	&= \frac{1}{z} + \sum_{n>0} \frac{2z}{z^2 - (n\pi)^2}
\end{align}

\newpage

\subsubsection{Verifica dell'espansione}

Dimostrare la validità della relazione
\begin{align}
	\cot(z) = \frac{1}{z} + \sum_{n>0} \frac{2z}{z^2 - (n\pi)^2}
\end{align}

Iniziamo dimostrando l'analiticità. Come abbiamo visto prima, dobbiamo prendere
\begin{align}
	f(z) = \cot(z) - \frac{1}{z}
\end{align}
Andiamo quindi a calcolarne il limite in 0
\begin{align}
	\limit{z}{0} \cot(z) \frac{1}{z} &= \limit{z}{0} \frac{\frac{z\cos(z)}{\sin(z)} -1}{z} = \nonumber\\
	&= \limit{z}{0} \frac{z\cos(z) - \sin(z)}{z\sin (z)} = \nonumber\\
	&= \limit{z}{0}\frac{\cos(z) \cdot [z - \tan (z)]}{z - \sin(z)} = \limit{z}{0} \frac{\cot(z)}{z} \cdot(z-\tan(z)) = 1\cdot 0 = 0
\end{align}
Verificata l'analiticità, procediamo a studiare poli e residui della funzione. 

Avremo dei poli semplici quando $\sin(z) = 0$, e qundi per $z_n=n\pi$, con $n=0$ \textbf{escluso}, in quanto vi si annulla la funzione. Ne segue che i residui saranno
\begin{align}
	Res(f,z_n) = \limit{z}{n\pi}  (z-n\pi) \cdot \frac{z\cot(z) -1}{z}
\end{align} 
Facendo il cambio di variaible $u=z-n\pi$ otteniamo
\begin{align}
	Res(f,z_n) = \limit{u}{0} u \cdot \frac{(u+n\pi)\cot(u+n\pi) -1}{u+n\pi} \label{eq3}
\end{align} 
Se però osserviamo un attimo il termine $\cot(u+n\pi)$ notiamo come lo si possa manipolare nel seguente modo
\begin{align}
	\cot(u+n\pi) &= \frac{\cos(u+n\pi)}{\sin(u+n\pi)} = \frac{\cos(u)\cdot(-1)^n}{\sin(u)\cdot(-1)^n} = \frac{\cos(u)}{\sin(u)} = \cot(u)
\end{align}
Possiamo quindi riscrivere la \ref{eq3} come
\begin{align}
	Res(f,z_n) &= \limit{u}{0} u \cdot \frac{(u+n\pi)\cot(u) -1}{u+n\pi} \simeq \limit{u}{0} u\cdot \frac{(u+n\pi)\cdot \frac{1}{u} -1}{u+n\pi} = \nonumber\\
	&= \left. \frac{u+n\pi -u}{u+n\pi} \right|_{u=0} = 1 \quad \forall n\in \N
\end{align} 

Possiamo quindi applicare lo sviluppo
\begin{align}
	f(z) &= \sum_{n\neq 0} \left( \frac{1}{n-n\pi} + \frac{1}{n\pi} \right) = \nonumber\\
		 &= \sum_{n=-\infty}^{-1} \left( \frac{1}{n-n\pi} + \frac{1}{n\pi} \right) + \sum_{n=1}^{+\infty} \left( \frac{1}{n-n\pi} + \frac{1}{n\pi} \right) = \nonumber \\
		 &= \sum_{n>0} \left[ \frac{1}{z-n\pi} + \frac{1}{n\pi} - \left( \frac{1}{-z-n\pi} + \frac{1}{n\pi} \right) \right] = \nonumber\\
		 &= 2z\sum_{n=1}^{+\infty} \frac{1}{z^2 - (n\pi)^2}
\end{align}
E la tesi è così dimostrata.

\newpage

\subsection{Esempio: $\sin^{-2}(z)$}

Utilizziamo l'integrale (di Cauchy? Chiedere)
\begin{align}
	I_n = \oint_{Q_n} \frac{d\zeta}{(\zeta -z) \sin^2(\zeta)}
\end{align}
Studiando i poli otteniamo
\begin{align}
	\zeta_0 &= z \quad &\text{polo semplice}\\
	\zeta_n &= n\pi \; n\in \N  \quad &\text{poli doppi}
\end{align}
Siccome $I_n \arrowlim{n}{+\infty} 0$ abbiamo che
\begin{align}
	0 &= \sum_{n\in \N} Res(f(\zeta), n\pi) + Res(f(\zeta), z) =\nonumber\\
	  &= \sum_{n\in \N} Res(f(\zeta), n\pi) + \frac{1}{\sin^2(z)}  \label{eq3}
\end{align}
Da cui otteniamo
\begin{align}
	  \frac{1}{\sin^2(\zeta)} &= -\sum_{n\in \N} Res(f(\zeta), n\pi)
\end{align}
Dato che
\begin{align}
	Res(f(\zeta), n\pi) &= \limit{\zeta}{n\pi} \frac{d}{d\zeta} \frac{(\zeta - n\pi)^2}{(\zeta - z) \sin^2(\zeta)} \firstpassage
	\double{u=\zeta - n\pi}{\zeta = u + n\pi} \nextpassage
	Res(f(\zeta), n\pi) &= \limit{u}{0} \frac{d}{du} \frac{u^2}{(u + n\pi - z) \sin^2(u + n\pi)} =  \limit{u}{0} \frac{d}{du} \frac{u^2}{(u + n\pi - z) (-1)^{2n} \sin^2(u)}= \nonumber\\
	&= \limit{u}{0} \frac{d}{du} \frac{u^2}{(u + n\pi - z) \left(u -\frac{u^3}{3!} + \dots\right)^2} = \left. \frac{d}{du} \frac{u^2}{(u + n\pi - z) \left(u -\frac{u^3}{3!} + \dots\right)^2} \right|_{u=0} = \nonumber \\
	&= \left. \frac{d}{du} \frac{1}{(u + n\pi - z) \left(1 -\frac{u^2}{3!} + \dots\right)^2} \right|_{u=0} = \left. \frac{d}{du} \frac{1}{(u + n\pi - z) \left(1 -\frac{u^2}{3!} + \dots\right)^2} \right|_{u=0} = \nonumber \\
	&= \left. \frac{d}{du} \frac{\left(1 +\frac{u^2}{3!} - \dots\right)^2}{(u + n\pi - z) } \right|_{u=0} = \left. \frac{d}{du} \frac{\left(1 +\frac{u^2}{3} + \dots\right)}{(u + n\pi - z) } \right|_{u=0} = \nonumber \\
	&= \left. \frac{\left( \frac{2}{3} u + \dots \right) (u + n\pi-z ) - \left(1 + \frac{u^2}{3}\right)}{(u+ n\pi -z)^2} \right|_{u=0} = -\frac{1}{(n\pi -z)^2}
\end{align}
La \ref{eq3} diventa
\begin{align}
	\frac{1}{\sin^2(z)} = \sum_{n\in \N} \frac{1}{(n\pi -z)^2}
\end{align}

\newpage

\section{Sommerfeld-Watson (Davide Bufalini)}

Utile strumento per il calcolo delle somme di serie. Sia una funzione $g(z)$ analitica ovunque tranne che in singolarità polari isolate $z_k$ e sia 
\begin{align}
	\limit{|z|}{+\infty} |zq(z) =0|
\end{align}
Allora valgono le seguenti relazioni
\begin{align}
	\sum_{n\in \Z} (-1)^n g(n) &= -\pi \sum_{k} Res\left( \frac{g(z)}{\sin(\pi z)}, z=z_k \right)\\
	\sum_{n\in \Z} g(n)  &= -\pi \sum_{k} Res\left( g(z)\cot(\pi z), z=z_k \right)
\end{align}

\subsection{Esempio 1}

Calcolare la somma della serie
\begin{align}
	S(a) = \sum_{n=1}^{+\infty} \frac{(-1)^n}{n^2 + a^2} \spacer a \in \R \backslash \{0\}
\end{align}
Iniziamo notando che possiamo scrivere
\begin{align}
	\sum_{n=1}^{+\infty} \frac{(-1)^n}{n^2 + a^2} = 2S(a) + \frac{1}{a^2}
\end{align}
Ponendo
\begin{align}
	g(z)=\frac{1}{n^2 + a^2}
\end{align}
Otteniamo
\begin{align}
	&z^2 + a^2= 0 \firstpassage
	&z_p= \pm ia	
\end{align}
E possiamo quindi applicare SW
\begin{align}
	&\limit{z}{\pm ia} \frac{z\pm ia}{z^2+a^2} \cdot \frac{1}{\sin(\pi z)} = \dots = -\frac{1}{2a\sinh(\pi a)} \firstpassage
	&\frac{(-1)^n}{n^2 + a^2} = -\pi Res \left(\frac{g(z)}{\sin(\pi z)}, z_p\right) = -\frac{\pi}{2a\sinh(\pi a)} \nextpassage
	&S(a) = -\frac{1}{2a^2} + \frac{\pi}{4a\sinh(\pi a)}
\end{align}


\newpage

\section{Funzioni speciali (Alessandro Marcelli)}

\subsection{Funzione Gamma di Eulero}
\begin{align}
	\Gamma(z) &= \int_{0}^{\infty} dt \; e^{-t} t^{z-1} = \dots =  \quad &\text{Integrale Euleriano di II tipo}\\
	&= \frac{1}{z} \prod_{k=1}^{\infty} \left( 1+ \frac{1}{k} \right)^z \left( 1+ \frac{z}{k} \right)^{-1}
\end{align}
Gode della proprietà	
\begin{align}
	&\Gamma(z) \Gamma(1-z) = \frac{\pi}{\sin(\pi z)} \firstpassage
	&\Gamma\left(\frac{1}{2}\right) \Gamma\left(\frac{1}{2}\right) = \pi \rightarrow \Gamma \left( \frac{1}{2} \right) = \sqrt{\pi} \\
	&\Gamma\left(\frac{1}{2} + \xi\right) \Gamma\left(\frac{1}{2} - \xi\right) = \frac{\pi}{\cos(\pi \xi)}\\
	&\sqrt{\pi} \Gamma(2z) = 2^{2z-1} \Gamma(z) \Gamma \left( z + \frac{1}{2} \right)\\
	&\Gamma(\overline{z}) = \overline{\Gamma(z)}
\end{align}
In campo complesso si può prolungare tramite il cammino di Hankel, e si arriva ad ottenere
\begin{align}
	\Gamma(n) = (n-1)!
\end{align}

\subsection{Funzione Digamma}
\begin{align}
	\psi(z) = \frac{\Gamma'(z)}{\Gamma(z)} = \frac{d}{dz} \ln \Gamma (z)
\end{align}
\subsection{Funzione beta di Eulero}
\begin{align}
	B(p,q) &= \int_{0}^{1} dt \; t^{p-1} (1-t)^{q-1} = \dots =  \quad &\text{Integrale Euleriano di I tipo}\\
	&= \frac{\Gamma(p)\Gamma(q)}{\Gamma(p+q)}
\end{align}

\subsection{Funzione zeta di Riemann generalizzata}

Sono stanco capo.




\section{Espansione di Weierstrass (Davide Bufalini)}
Presa una $f(z)$ con infiniti zeri $z_j$ numerabili di ordine $\alpha_j$ che si accumulano all'infinito (e mai allo zero???) allora vale lo sviluppo
\begin{align}
	f(z) =f(0) e^{z\frac{f'(0)}{f(0)}} \prod_{n=1}^{+\infty}\left( 1 - \frac{z}{z_n} \right)^{\alpha_n} e^{\alpha_n \frac{z}{z_n}}
\end{align}

\newpage

\section{Espansioni asintotiche (Davide Bufalini, Alessandro Marcelli, Paolo Proia)}

\subsection{Spunti di teoria (Alessandro Marcelli, Paolo Proia)}

\subsubsection{Definizione molto generale (Alessandro Marcelli)}

Lo sviluppo in serie può essere esteso anche al caso in cui la funzione oggetto di studio $f(z)$ non sia analitica in $z_0$, trovando uno sviluppo $f(z) \sim \sum_{n=1}^{\infty} a_n f_n(z)$ che approssima "bene" il comportamento della funzione in un intorno di $z_0$ \cite{MMF}, ovvero che valgono le due seguenti proprietà:
\begin{enumerate}
	\item il resto della serie $R_N(z) = f(z) - S_N(z)$ è tale che $\limit{z}{z_0} \frac{|R_N(z)|}{|f_N(z)|} = 0 \quad \forall N\in \N$
	\item la serie fornisce una buona approssimazione di $f(z)$ solo in determinate regioni del piano complesso
\end{enumerate}
L'approssimazione che si va a trovare sarà tanto più buona tanto più grande è $x$, e si lavora quindi per $x \to + \infty$. Nel caso ci si trovi nel caso $x \to -\infty$ bisogna fare un cambio di variabile.

Il contributo maggiore sarà dato dai primi termini dell'espansione, motivo per cui di solito ci si ferma al secondo termine.

Nel nostro caso ci interessiamo della famiglia di funzioni integrali del tipo
\begin{align}
	F(x) = \int_{a}^{b} dt \; e^{x\phi(t)}g(t) \label{lol0}
\end{align}

\subsubsection{Termine Leading: Metodo di Laplace (Alessandro Marcelli)}

Se interessa solo il primo termine dell'espansione, detto \textbf{dominante}, basta sviluppare al secondo ordine $\phi(t)$ e $g(t)$ nel punto di massimo della prima, risolvere l'integrale con questa approssimazione e prendere il primo termine, come si vede nell'esempio svolto del paragrafo \ref{es1}.

\subsubsection{Metodo risolutivo generalizzato (Paolo Proia)}

Qualora servano termini successivi (di solito ci si ferma al termine Next-to-Leading) la situazione si complica, in quanto non basta sviluppare al secondo ordine e bon, in quanto ci sono casi in cui uno dei due termini può essere nullo (funzioni trigonometriche my beloved). Il procedimento si snoda quindi:
\begin{enumerate}
	\item Trovare massimo per \(\phi(t)\) (chiamiamo \(t_0\))
	\item Espandere \(g(t)\) intorno a \(t_0\) per 2 o 3 termini
	\item Prendere nota dell'esponente dell'ultimo termine (definiamo come  \(n\)) 
	\item Parametrizzare come $\phi(t) - \phi(t_0)=-\tau^2 \label{lol1}$
	ed espandere \(\phi(t)\) intorno a $t_0$ fino ad ordine \(n\)
	\item Trovare cambio di variabile da \(t\) a \(\tau\), imponendo $t - t_0 = h(\tau) = \sum_i^{n+1} c_i \tau^i$
	\item Sostituiamo nella \refeq{lol1} tenendo i termini fino a \(\tau^{n+1}\)
	\item Imporre l'uguaglianza tra i coefficenti con la stessa potenza di \(\tau\) a dx e sx, trovando così \(h(\tau)\)
	\item Essendo \(h(\tau) = t + cost. \) abbiamo che \(\dot{h}d\tau = dt\) e abbiamo così terminato il cambio di variabile
	\item Sostituiamo nella \ref{lol0}
	\begin{align}
		&\phi(t) = \phi(t_0) - \tau^2\\
		&dt=\dot{h}d\tau\\
		&g(t) \to g(h(\tau)) \label{lol2}
	\end{align}
	Nella \ref{lol2} usiamo direttamente l'espansione in serie del punto 2
	\item Svolgiamo il prodotto \(\dot{h} \cdot g(h(\tau))\), sostituiamo nell'integrale e teniamo i primi $n$ termini il cui integrale è non nullo (nel caso Next-to-Leading i primi 2)
\end{enumerate}

\textbf{NOTA BENE:} qualora una delle due funzioni sia una trigonometrica elevata a potenza conviene sviluppare solo la trigonometrica e mettere a potenza lo sviluppo, e fare altri magheggi
quali utilizzare le relazioni trigonometriche per fare il meno conti possibile.

\subsection{Esempio 1: metodo generalizzato (Paolo Proia)}

Sia la seguente funzione integrale
\begin{align}
	I(t)    &= \int_{1}^{+\infty} dt \; e^{-x \frac{t^2}{2}}t^x \frac{1}{t^2 + 1} = \int_{1}^{+\infty} dt \; e^{-x \frac{t^2}{2}} e^{x\ln(t)} \frac{1}{t^2 + 1} = \int_{1}^{+\infty} dt \; e^{x\left(\ln(t) -\frac{t^2}{2}\right)} \frac{1}{t^2 + 1}\\
	\phi(t) &= \ln(t) - \frac{t^2}{2} \\
	g(t)    &= \frac{1}{t^2 + 1}
\end{align}

\begin{enumerate}
	\item Trovare massimo per \(\phi(y)\) (chiamiamo \(y_0\))
	\begin{align}
		\dot{\phi}(t) &=\frac{1}{t} - t = \frac{1- t^2}{t} = 0 \to t = \pm 1\\
		\ddot{\phi}(t) &= -\frac{1}{t^2} - 1 = -\frac{t^2+1}{t^2}  \\
		\dddot{\phi}(t) &= \frac{2}{t^3}
	\end{align}
	Escoludiamo il risulato negativo, siccome $\phi(t)$ non è definito in quel punto. Abbiamo quindi
	\begin{align}
		t_0=+1 \quad ; \quad\phi_0 = -\frac{1}{2}
	\end{align}
	\item Espandiamo \(g(t)\) intorno a \(t_0\)
	\begin{align}
		g(t) &\simeq g(t_0) + g(t_0)\cdot(t-t_0) = \frac{1}{2} - \frac{1}{2} (t-1) = 1-\frac{t}{2} 
	\end{align}
	\item Notiamo come ci siamo fermati per \(n=1\)
	\item Parametrizziamo
	\begin{align}
		&\phi(t)= -\frac{1}{2} + 0\cdot(t-1) - \frac{1}{2}\cdot 2 \cdot(t-1)^2  + \frac{1}{3!} \cdot 2 (t-t_0)^3= -\frac{1}{2} - (t-1)^2 + \frac{1}{3} (t-t_0)^3= -\tau^2
	\end{align}
	\item Troviamo cambio di variabile da \(t\) a \(\tau\), fermandoci al terzo ordine, essendo lo sviluppo in t al 3o ordine
	\begin{align}
		&t-t_0 = t-1 = c_1 \tau + c_2 \tau^2 + c_3 \tau^3 + O(\tau^4)\\
		&\phi(t) = - \frac{1}{2} - (c_1 \tau + c_2 \tau^2 + c_3 \tau^3)^2+ \frac{1}{3} (c_1 \tau + c_2 \tau^2 + c_3 \tau^3)^3
	\end{align}
	\item Andiamo quindi a sostituire
	\begin{align}
		\phi(t) + \frac{1}{2} &=  - (c_1 \tau + c_2 \tau^2 + c_3 \tau^3)^2+ \frac{1}{3} (c_1 \tau + c_2 \tau^2 + c_3 \tau^3)^3 = \dots = \nonumber\\
		&= -c_1^2 \tau^2 - 2c_1c_2 \tau^3 + \frac{1}{3}c_1^3 \tau^3
	\end{align}
	\item Andiamo ora a confrontare
	\begin{align}
		\phi(t) + \frac{1}{2} &= -\tau^2\\
		\phi(t) + \frac{1}{2} &= -c_1^2 \tau^2 - 2c_1c_2 \tau^3 + \frac{1}{3}c_1^3 \tau^3
	\end{align}
	Ne segue che
	\begin{align}
		&-c_1^2 \tau^2 = -\tau^2 &\to c_1 = 1\\
		&-2c_1c_2 \tau^3 + \frac{1}{3}c_1^3 \tau^3=0 &\to -2c_2 + \frac{1}{3} = 0 \to c_2 = \frac{1}{6}
	\end{align}
	Abbiamo così trovato
	\begin{align}
		t-1 = h(\tau) = \tau + \frac{1}{6}\tau^2
	\end{align}
	\item Sostuiamo e troviamo
	\begin{align}
		t-1 = \tau + \frac{1}{6}\tau^2 \to d(t-1) = dt = \left( 1 + \frac{1}{3} \tau \right) d\tau
	\end{align}
	\item Sostituiamo infine nella nostra funzione integrale
	\begin{align}
		I(x) &= \int_{0}^{+\infty} d\tau \; e^{-x \left( \tau^2 + \frac{1}{\tau} \right)} \left( 1 + \frac{1}{3} \tau \right) \left(\frac{1}{2} - \frac{\tau}{2}\right) = \nonumber\\
		&= \int_{0}^{+\infty} d\tau \; e^{-x \left( \tau^2 + \frac{1}{\tau} \right)} \left( \frac{1}{2}- \frac{1}{2}\tau + \frac{1}{6} \tau - \frac{1}{6} \tau^2 \right) = \nonumber\\
		&= \int_{0}^{+\infty} d\tau \; e^{-x \left( \tau^2 + \frac{1}{\tau} \right)} \left( \frac{1}{2} - \frac{1}{3} \tau \right) = \nonumber\\
		&= e^{-\frac{x}{2}}\int_{0}^{+\infty} d\tau \; e^{-x\tau^2} \left( \frac{1}{2} - \frac{1}{3} \tau \right)
	\end{align}
	\item abbiamo quindi
	\begin{align}
		I(x) &= e^{-\frac{x}{2}} \left( \frac{1}{2} \int_{0}^{+\infty} d\tau \; e^{-x\tau^2} - \frac{1}{3} \int_{0}^{+\infty} d\tau \; e^{-x\tau^2}\tau \right)
	\end{align}
\end{enumerate}

\newpage

\subsection{Esempio 2 (Davide Bufalini) \label{es1}}

Calcolare per $x\rightarrow +\infty$ i primi due termini dell'espansione per la funzione integrale
\begin{align}
	F(x) = \int_{0}^{\frac{\pi}{2}} dt \; e^{x\cos(t)}\sin^3(t)
\end{align}

Partiamo dal \textbf{termine Leading}. Siccome abbiamo
\begin{align}
	\phi(t) &= \cos(t) \\
	g(t) &= \sin^3(t)
\end{align}
Andiamo a cercarne il punto di massimo, che troviamo in
\begin{align}
	&\frac{d \phi(t)}{dt} = -\sin(t) \firstpassage
	&\sin(t) = 0 \nextpassage
	&t_0=0
\end{align}
Andiamo quindi a sviluppare
\begin{align}
	\phi(t) &\simeq 1- \frac{t^2}{2} + o(t^3)  \\
	g(t) &\simeq \left( t - \frac{t^3}{6} \right)^3 \simeq t^3 \quad \text{termine leading}
\end{align}
E otteniamo
\begin{align}
	F(x) &\simeq \int_{0}^{\frac{\pi}{2}} dt \; e^{x\left(1- \frac{t^2}{2}\right)}t^3 = \nonumber\\
	&= \int_{0}^{+\infty} dt \; e^x e^{-x\frac{t^2}{2}}t^3
\end{align}
Applicando la sostituzione
\begin{align}
	&y= \frac{t^2}{2} \quad \leftrightarrow \quad t^2 = 2y \\
	&dy = tdt
\end{align}
Possiamo arrivare ad ottenere il termine cercato
\begin{align}
	F(x) &= \int_{0}^{\frac{\pi}{2}} dy \; e^x e^{-xy}2y = \nonumber\\
		 &= 2e^x\int_{0}^{+\infty} dy \;e^{-xy}y = \nonumber\\
		 &= 2e^x \frac{\Gamma(2)}{x^2} = \frac{2e^x}{x^2}
\end{align}

Per quanto riguarda il termine \textbf{Sub-leading}, andiamo a parametrizzare la $\phi(t)$ nel seguente modo
\begin{align}
	&\phi(t) - \phi(t_0) = -\tau^2 \firstpassage
	&\cos(t) - 1 = -\tau^2 \nextpassage
	&t= \arccos(i-\tau^2) \firstpassage
	&dt = -\frac{d\tau}{\sqrt{1-(1-\tau^2)^2}}\cdot(-2\tau) = \frac{2\tau d\tau}{\sqrt{2\tau^2-\tau^4}} = \frac{2d\tau}{\sqrt{2-\tau^2}}
\end{align}
Andiamo così a trovare i primi due termini dell'espansione
\begin{align}
	F(x) &\simeq \int_{0}^{+\infty} \frac{2d\tau}{\sqrt{2-\tau^2}} \; e^{x\left(1- \tau^2\right)}(\sqrt{1-cos^2(t)})^3 = \nonumber\\
	&= \int_{0}^{+\infty} \frac{2d\tau}{\sqrt{2-\tau^2}} \; e^{x\left(1- \tau^2\right)}(\sqrt{1-(1-\tau^2)^2})^3 = \nonumber\\
	&=\int_{0}^{+\infty} \frac{2d\tau}{\sqrt{2-\tau^2}} \; e^{x\left(1- \tau^2\right)}(\sqrt{2\tau^2-t^4})^3 = \nonumber\\
	&=\int_{0}^{+\infty} \frac{2d\tau}{\sqrt{2-\tau^2}} \; e^{x\left(1- \tau^2\right)}(\tau\sqrt{2-\tau^2})^3 = \nonumber\\
	&=\int_{0}^{+\infty} 2d\tau \; e^{x\left(1- \tau^2\right)}\tau^3 (2-\tau^2) = \nonumber\\
	&=2e^x \left[ 2\int_{0}^{+\infty}d\tau \; e^{-x\tau^2} \tau^3 - \int_{0}^{+\infty} d\tau\; e^{-x\tau^2}\tau^5 \right] = \nonumber\\
	&= 2e^x\left[ 2\frac{\Gamma(2)}{x^2} - \frac{1}{2} \frac{\Gamma(3)}{x^3} \right] = \nonumber \\
	&= \frac{2e^x}{x^2} - \frac{2}{2} \frac{e^x 2!}{x^3} =\nonumber\\
	&= \frac{2e^x}{x^2} -\frac{2e^x}{x^3}
\end{align}



\subsection{Esempio 3 (Davide Bufalini)}

Calcoliamo i termini \textbf{Leading} e \textbf{Subleading} per la funzione
\begin{align}
	I(x) = \int_\R dt \; e^{-x \cosh(t)}
\end{align}
Iniziamo individuando le  nostre funzioni
\begin{align}
	\phi(t) &= -\cosh(t)\\
	g(t) &= 1 \firstpassage
	\frac{d\phi(t)}{dt} &= \sinh(t) = 0 \quad \rightarrow \quad t_0 = 0
\end{align}

Inziamo col termine \textbf{Leading}. Sviluppando intorno a $t_0$ otteniamo
\begin{align}
	&\phi(t) \simeq -\left( 1 + \frac{t^2}{2} \right)\firstpassage
	&I(x) = \int_\R dt e^{-x}e^{-x \frac{t^2}{2}} = e^{-x} \frac{2\pi}{x}	
\end{align}

Per quanto riguarda il termine \textbf{Subleading} andiamo a parametrizzare
\begin{align}
	&\phi(t) - \phi(t_0) = -\tau^2 \firstpassage
	& \cosh(t) -(-1) = -\tau^2
\end{align}
Siccome abbiamo
\begin{align}
	dt = \left( \frac{dt}{d\tau} \right) d\tau
\end{align}
Dobbiamo trovare una espressione per $t(\tau)$ per il calcolo ello Jacobiamo.

Siccome abbiamo $t_0 = 0$ andiamo a sviluppare per $t$ piccoli
\begin{align}
	&\cosh(t) = 1 + \frac{t^2}{2} + \frac{t^4}{24} + \dots \firstpassage
	&\cancel{1} + \frac{t^2}{2} + \frac{t^4}{24} + \dots = \cancel{1} + \tau^2 \nextpassage
	&\frac{t^2}{2} \left( 1 + \frac{t^2}{12} \right) = \tau^2 \firstpassage
	&t^2 = \frac{2\tau^2}{1 + \frac{t^2}{12}} = 2\tau^2 \left( 1 - \frac{t^2}{12} + \left( \frac{t^2}{12} \right)^2 + \dots \right) \simeq 2\tau^2 - 2\tau^2 \frac{t^2}{12} \firstpassage
	&t^2 \left( 1 + \frac{\tau^2}{6} \right) = 2\tau^2 \nextpassage
	&t^2=\frac{2\tau^2}{1 + \frac{\tau^2}{6}}\nextpassage
	&t=\frac{\sqrt{2}\tau}{\sqrt{1 + \frac{\tau^2}{6}}} \simeq \sqrt{2}\tau \left( 1 - \frac{1}{2} \frac{\tau^2}{6} \right) = \sqrt{2}\tau \left( 1 - \frac{\tau^2}{12} \right)
\end{align}
Dunque otteniamo
\begin{align}
	t &= \sqrt{2} \tau - \frac{\sqrt{2}}{12}\tau^3 \firstpassage
	dt &= \left( \sqrt{2} - \frac{\sqrt{2}}{4}\tau^2 \right)d\tau 
\end{align}

Possiamo quindi finalmente sostituire
\begin{align}
	I(x) &= \int_\R dt \; e^{-x} e^{-\frac{t^2}{2} x} = \nonumber \\
		 &= e^{-x} \int_\R d\tau \; e^{-x\tau^2}  \left( \sqrt{2} - \frac{\sqrt{2}}{4}\tau^2 \right) = \nonumber \\
		 &= e^{-x}\sqrt{2} \left[ \int_\R d\tau \; e^{-x\tau^2}   -\int_\R d\tau \;e^{-x\tau^2}  \frac{1}{4}\tau^2  \right]= \nonumber \\
		 &= e^{-x}\sqrt{2} \left[ \sqrt{\frac{\pi}{x}}   - \frac{1}{4}\cdot\frac{1}{x} \cdot \frac{1}{2} \cdot \sqrt{\frac{\pi}{x}}  \right]= \nonumber \\
		 &=e^{-x}\sqrt{\frac{2\pi}{x}} - \frac{1}{8x}e^{-x}\sqrt{\frac{2\pi}{x}}= \nonumber \\
		 &=e^{-x}\sqrt{\frac{2\pi}{x}}\left[1 - \frac{1}{8x}\right]
\end{align}

\newpage

\subsection{Esempio 4 (Davide Bufalini)}

Data la funzione
\begin{align}
	\int_{\frac{\pi}{2}}^{\frac{3}{4}\pi} dt \; e^{-\frac{x}{\sin(t)}} (1 + \cot (t))
\end{align}
se ne calcoli l'espansione.

Iniziamo trovando il punto di sella
\begin{align}
	&\phi(t) = -\frac{1}{\sin(t)}\firstpassage
	&\frac{d\phi}{dt} = \frac{\cos (t)}{\sin^2 (t)} = 0 \firstpassage
	& t_0 = \frac{\pi}{2}
\end{align}
Siccome il punto cade su uno degli estremi, ci basta considerare $t(\tau) \sim o(\tau^2)$
Parametrizziamo
\begin{align}
	&\phi(t) - \phi(t_0) = -\tau^2\firstpassage
	&-\frac{1}{\sin (t)}-\left(-\frac{1}{\sin \left( \frac{\pi}{2}\right)}\right) = -\tau^2 \nextpassage
	&-\frac{1}{\sin (t)} = 1 +\tau^2\nextpassage
	&\sin (t) = -\frac{1}{1+\tau^2} = 1 - \tau^2 + \tau^4 + \dots \nextpassage
	&\frac{1}{2} \left( t -\frac{\pi}{2} \right)^2 = \tau^2 + o(\tau^3) \label{kekw}
\end{align}
Espandendo $\tau$ si ha che
\begin{align}
	 t -\frac{\pi}{2} = c_1 \tau + c_2 \tau^2 + c_3 \tau^3 + \dots \label{kukw}
\end{align}
Sostituendo nella \ref{kekw}
\begin{align}
&\frac{1}{2} \left( c_1 \tau + c_2 \tau^2 \right)^2 = \tau^2 \firstpassage
&c_1^2 \tau^2 + 2c_1 c_2 \tau^3 + c_2^2 \tau^4 = 2 \tau^2\nextpassage
&\double{c_1 = &\sqrt{2}}{c_2 = &0}
\end{align}
Sostituendo nella \ref{kukw} otteniamo
\begin{align}
	&t -\frac{\pi}{2} = \sqrt{2}\tau + o(\tau^3)	\firstpassage
	&dt = \sqrt{2}d\tau
\end{align}
Sviluppiamo ora $g(t)$ intorno a $t_0$
\begin{align}
	g(t) = 1 + \cot(t) = 1 - \left( t -\frac{\pi}{2} \right) \sim 1 - \sqrt{2}\tau +o(\tau^3)
\end{align}
Possiamo quindi riscrivere la nostra funzione integrale come
\begin{align}
	I(x) &= \int_{0}^{+\infty} \sqrt{2} d\tau \; e^{-x(1+\tau^2)}(1 - \sqrt{2}\tau) = \nonumber\\
	&= \sqrt{2}e^{-x} \int_{0}^{+\infty}  d\tau \; e^{-\tau^2}(1 - \sqrt{2}\tau) = \nonumber\\
	&=e^{-x} \left[ \sqrt{\frac{\pi}{2x}} - \frac{1}{x} \right]
\end{align}


\newpage

\subsection{Metodo del punto di sella (Davide Bufalini)}

Abbiamo finora visto integrali sulla retta reale. Qualora si abbia una generica curva $\gamma \in \C$ il procedimento si complica leggermente. 

Presa una funzione integrale della forma
\begin{align}
	I(x) = \int_\gamma dt \; e^{x\phi(t)} g(t)
\end{align}
Qualora $\phi(t)$ e $g(t)$ siano analitiche in un dominio $\mathit{D}$ e si abbia $\gamma \subseteq \mathit{D}$ abbiamo un problema. Infatti, riscrivendo la $\phi(t)$ nella seguente forma
\begin{align}
	\phi(t) &= u(t) + iv(t)
\end{align}
All'interno del dominio non è detto che $v(t)$ sia costante, il che ci impedisce di restringerci al caso reale. Inoltre, per $t \in \gamma$ il termine $e^{iv(t)}$ oscilla fortemente, portanto la $I(x)$ ad essere mediata a 0. 

Cerchiamo un modo per rendere costante la $v(t)$ e rendere massima la $u(t)$, in modo da rispettare il comportamento asintotico.
Il procedimento si snoda nei seguenti passaggi
\begin{enumerate}
	\item Riscriviamo la $\phi(t)$ come
	\begin{align}
		\phi(t) = \phi(x,y) = u(x,y) + iv(x,y)
	\end{align}
	Possiamo farlo in più modi
	\begin{enumerate}
		\item Coordinate cartesiane, sostituendo $t = x + iy$
		\item Coordinate polari, sostituendo $t = \sqrt{x^2 + y^2}\cdot e^{i \arctan \left(\frac{x}{y}\right)}$
	\end{enumerate}
	\item Cerchiamo il punto di sella $t_0$ di $\phi(t)$ tale
	\begin{align}
		\phi(t_0) &= 0 \\
		\ddot{\phi}(t_0) &: \ddot{u}(t_0) <0 
	\end{align}
	\item Cerchiamo quindi i percorsi a fase costante passanti per $t_0$, ovvero ponendo $v(t)= v(t_0)$, da cui segue $Im(\phi(t)) m= Im(\phi(t_0))$. Una volta trovati questi cammini, vediamo se è possibile deformare la curva $\gamma$ su un cammino a fase costante $\gamma^*$ dove $\ddot{u}(t_0) <0$  
	\item Una volta trovato questo cammino, risolviamo con i metodi noti per il caso reale
	
\end{enumerate}

\subsection{Esempio 1}

Calcoliamo il termine leading per l'espansione asintotica  della funzione integrale
\begin{align}
	F(x) &= \int_\R dy \; e^{-x[2y^2 -(y-i)^2]} \frac{1}{y^2 + 2}\\
	g(y) &=  \frac{1}{y^2 + 2}\\
	\phi(y) &= 2y^2 -(y-i)^2
\end{align}

Iniziamo trovando il punto di sella di $\phi(y)$, e calcoliamo preliminarmente
\begin{align}
	\phi(y) &=  (y-i)^2 -2y^2 = y^2 - 2iy -1 -2y^2 = -y^2 - 2iy -1 \\
	\dot{\phi}(y) &= 2y + 2i\\
	\ddot{\phi}(y) &= 2
\end{align}
Troviamo quindi che
\begin{align}
	&2y + 2i = 0 \firstpassage
	&y_0=-i\firstpassage
	&y_0= (0,-1)
\end{align}

Passiamo in coordinate cartesiane sostituendo $y = a+ib$
\begin{align}
	\phi(y)) &= \phi(a,b) = -(a+ib)^2 -2i(a+ib) -1 = \nonumber\\
	&= -a^2 - 2iab + b^2 - 2ia + 2b -1 =   \nonumber\\
	&= (b^2 -a^2 + 2b -1) -2ia(b + 1)
\end{align}
Da cui abbiamo che
\begin{align}
	u(a,b) &= b^2 -a^2 + 2b -1\\
	v(a,b) &= -2a(b + 1)
\end{align}
Cerchiamo ora i cammini a fase costante passanti per $y_0$. Sostituendo troviamo che
\begin{align}
	&v(0,-1) = -2\cdot 0 \cdot (0 + 1) =0\\
	&a(b + 1) = 0 \to \double{a=0 \quad &\text{asse immaginario}}{b=-1 \quad &\text{asse reale traslato di $-i$}}
\end{align}
Ma quale dei due cammini è quello che serve a noi? Andiamo a sostituire nell'espressione di $u(a,b)$:
\begin{enumerate}
	\item Sostituiamo $a=0$
	\begin{align}
		u(0,b) &= b^2 + 2b-1 \\
		\left. \frac{d u}{d b} \right|_{a=0} &= 2b + 2 \\ 
		\left. \frac{d^2 u}{d b^2} \right|_{a=0} &= 2 \geq 0  
	\end{align}
	Siccome non abbiamo $\ddot u <0$ questo cammino non va bene.
	\item Sostituiamo $b=-1$
	\begin{align}
		u(a,-1) &= 1 -a^2 -2-1 = -a^2 -2 \\
		\left. \frac{d u}{d a} \right|_{b=-1} &= -2a\\ 
		\left. \frac{d^2 u}{d a^2} \right|_{b=-1} &= -2 \leq 0  
	\end{align}
\end{enumerate}

Il cammino giusto è quindi 
\begin{align}
	\gamma^*= \left\{ (a,b) \in \C \taleche (a,b) = (a,-1) \right\}
\end{align}
I cui punti sono parametrizzati come
\begin{align}
	y= a - i \spacer a \in \R
\end{align}
E da cui segue che
\begin{align}
	Re(\phi) &= u(a,b) = -a^2 - 2 \\
	Im (\phi) &=v(a,b) = 0
\end{align}
Andiamo quindi a sostituire nella nostra funzione integrale ed otteniamo
\begin{align}
	F(x) &= \int_\R da \; e^{x[-2-a^2]} \frac{1}{a^2 -2ia -1 + 2} =\nonumber\\
	&= \int_\R da \; e^{x[-2-a^2]} \frac{1}{a^2 -2ia +1} 
\end{align}

Siccome stiamo su sviluppando sul punto di sella $(0,-1)$ il denominatore può essre approssimato a $1$, e abbiamo quindi
\begin{align}
	F(x) &= e^{-2x}\int_\R da \; e^{-xa^2} = e^{-2x} \sqrt{\frac{\pi}{x}}
\end{align}

\newpage

\subsection{Esempio 2}

Calcoliamo i termini Leading e Subleading per lo sviluppo asintotico della seguente funzione integrale
\begin{align}
	F(x) &= \int_{-i\infty}^{+i\infty} dx \; e^{ix(z^2-2z)} \frac{1}{\sqrt{1+z}}\\
	g(z) &=\frac{1}{\sqrt{1+z}}\\
	\phi(z) &= i(z^2-2z)
\end{align}
Andiamo a calcolare le derivate di $\phi(z)$ e a trovare il punto di sella
\begin{align}
	\phi(z) &= i(z^2-2z)\\
	\dot{\phi}(z) &= i(2z-2) \to z_0=1\\
	\ddot{\phi}(z) &= 2i
\end{align}
Passiamo in coordinate cartesiane ponendo $z = x+iy$ e otteniamo
\begin{align}
	z_0 &= x_0 + iy_0 = 1 \to z_0=(1,0)\\
	\phi(x,y) &= i[( x+iy )^2 -2x -2iy] = i( x^2 + 2ixy -y^2  -2x -2iy) = \nonumber\\
	&= i[( x^2 -y^2 -2x ) + 2i (xy -y)]  = \nonumber\\
	&=2y(1-x) + i( x^2 -y^2 -2x )
\end{align}
Da cui quindi otteniamo
\begin{align}
	\phi(x,y) &= u(x,y) + iv(x,y)\\
	u(x,y) &= 2y(1-x) \to u(0,1) = 2\\
	v(x,y) &= x^2 -y^2 -2x \to v(0,1) = -1
\end{align}
Vogliamo quindi trovare i percorsi tali che $Im(\phi(z)) = -1$. Per fare ciò poniamo
\begin{align}
	&x^2 -y^2 -2x = -1\firstpassage
	&x^2 -y^2 -2x +1 =0\nextpassage
	&(x-1)^2 = y^2\nextpassage
	&y=\pm(x-1)
\end{align}
Le nostre curve sono quindi due rette. Cerchiamo ora quale fa al caso nostro
\begin{align}
	y_1=x-1 \to u_1(x) &= -2(1-x)^2 \to \dot{u}_1(x) = -4(1-x)\\
	y_2=1-x \to u_2(x) &= +2(1-x)^2 \to \dot{u}_2(x) = +4(1-x)
\end{align}

Studiando l'andamento delle derivate notiamo come $u_1(x)$ abbia un massimo in $x=1$,e che quindi è la nostra curva d'interesse.

Andiamo ora a parametrizzare. Siccome i punti della nostra curva sono ruotati di $45°$ rispetto all'asse reale, scriviamo
\begin{align}
	&z = 1+ e^{i\frac{\pi}{4}} \tau\firstpassage
	&dz = e^{i\frac{\pi}{4}} d\tau
\end{align}	
Possiamo quindi riscrivere
\begin{align}	
	\phi(\tau) &= i[(1+ e^{i\frac{\pi}{4}} \tau)^2-2(1+ e^{i\frac{\pi}{4}} \tau)] =\nonumber\\
	&= i[ 1 + \cancel{2e^{i\frac{\pi}{4}}\tau} + e^{i\frac{\pi}{2}}\tau^2 - 2 -\cancel{2e^{i\frac{\pi}{4}}\tau} ]=\nonumber\\
	&= i (e^{i\frac{\pi}{2}}\tau^2 -1)=\nonumber\\
	&= i (-i\tau^2 -1) = -\tau^2 -i
\end{align}

\textbf{Nota bene:} abbiamo bisogno del termine subleading, ma questo non compare dalla formula del $d\tau$ che si ferma prima di $\tau^2$. Dobbiamo quindi sviluppare $g(z)$ intorno a $z_0$
\begin{align}
	g(z) = \left. \frac{1}{\sqrt{1+z}} \right|_{z=1 + e^{1\frac{\pi}{2}}} = \frac{1}{\sqrt{2}} \frac{1}{\sqrt{1 + \frac{e^{1\frac{\pi}{2}}}{2}}} \sim 1 - \cancel{\frac{e^{i\frac{\pi}{4}}}{4}\tau} + \frac{3}{8} \frac{e^{i\frac{\pi}{2}}}{4} \tau^2
\end{align}
Possiamo quindi riscrivere la nostra funzione integrale come
\begin{align}
	I(x) &= \int_\R d\tau; \frac{e^{x(-i-\tau^2)}}{\sqrt{2}} \left( 1  + \frac{3}{32} i \tau^2 \right)e^{i\frac{\pi}{4}} = \nonumber \\
	&= \frac{e^{-i\left( x-\frac{\pi}{4} \right)}}{\sqrt{2}}\left[ \int_\R d\tau  \; e^{-x\tau^2}  + \frac{3i}{32} \int_\R d\tau \; e^{-x\tau^2} \tau^2\right] = \nonumber \\
	&= \frac{e^{-i\left( x-\frac{\pi}{4} \right)}}{\sqrt{2}}\left[ \sqrt{\frac{\pi}{x}}   + \frac{3i}{32} \sqrt{\frac{\pi}{x^3}} \right]
\end{align}
\chapter{Operatori Infinito-Dimensionali}

\section{Spunti di teoria generale (Alessandro Marcelli)}

\subsection{Nozioni generali sugli spazi}

\subsubsection{Spazi vettoriali}

Si definisce uno \textbf{spazio vettoriale} $V$ come un insieme su un campo $K$ (di solito $\R$ o $\C$) dove sono definite le operazioni di
\begin{enumerate}
	\item \textbf{Somma vettoriale}
	\begin{align}
		V \times V &\rightarrow V\\
		(v_1,v_2)&\rightarrow v_1 + v_2
	\end{align}
	\item \textbf{Prodotto per uno scalare}
	\begin{align}
	K \times V &\rightarrow V\\
	(\lambda,v)&\rightarrow \lambda v
\end{align}
\end{enumerate}

Preso $A$ sottoinsieme di $V$, ne definiamo i \textbf{punti di aderenza} come quei punti tali che, preso un loro intorno, contengono almeno un punto di $A$. \newline
In particolare, preso un punto di aderenza $z_0$, se nel suo intorno contiene almeno un punto $z_1 \neq z_0$ appartenente ad $A$, allora $z_0$ si dice \textbf{punto di accumulazione}.

Possiamo quindi dividere i punti di aderenza in due categorie
\begin{enumerate}
	\item tutti i punti di A
	\item  i punti di accumulazione esterni ad $A$
\end{enumerate}

Possiamo ora definire la \textbf{chiusura} $\overline{A}$ come l'insieme dei punti di aderenza di $A$. 

Un insieme $A$ si dice \textbf{denso} in $B$ se $B \subseteq \overline{A}$. In particolare se $\overline{A} = V$ l'insieme $A$ si dice \textbf{ovunque denso}.



\newpage
\subsubsection{Spazi metrici, normati ed euclidei}

Dato uno spazio  vettoriale V e, presi due elementi dello spazio $x$ e $y$, definita la funzione \textbf{distanza} $d(x,y)$ tale che
\begin{align}
	d(x,y) &\geq 0 \spacer d(x,y) = 0 \leftrightarrow x=y\\
	d(x,y) & = d(y,x)\\
	d(x,y) &\leq d(x,z) + d(y,z \spacer \forall x,y,z \in M) 
\end{align}

Si definisce \textbf{Spazio Metrico} la coppia $(V,d)$.

Preso di nuovo lo spazio vettoriale $V$ e presi $v,w\in V$ e $\alpha \in K$ si può definire la \textbf{norma} $||\cdot|| \taleche V \rightarrow \R$ come
\begin{align}
	||v|| &\geq 0 \spacer ||v|| = 0 \leftrightarrow v= 0\\
	||\alpha v || &= |\alpha|\cdot ||v||\\
	||v + w|| &\leq ||v|| + ||w||
\end{align}

La coppia $(V, ||\cdot ||)$ si chiama \textbf{spazio normato}. 


\textbf{Importante:} tutti gli spazi normati sono anche spazi metrici, con $d(x,y) = ||x-y||$, ma non tutti gli spazi metrici sono normati.

In uno spazio metrico abbiamo che un punto $x_a\in M$ sarà
\begin{enumerate}
	\item di \textbf{aderenza} se e solo se $\exists \{x_n\} \in M \taleche \limit{n}{+\infty} x_n = x_a$
	\item di \textbf{accumulazione} se e solo se $\exists \{x_n\neq x_a\} \in M \taleche \limit{n}{+\infty} x_n = x_a$
\end{enumerate}

Definiamo ora una \textbf{successione di Cauchy} come tale se
\begin{align}
	\forall \epsilon > 0 \; \exists n_\epsilon \taleche d(x_m,x_n) <\epsilon \quad \forall m,n > n_\epsilon
\end{align}
Possiamo quindi dire che uno spazio metrico $V$ si dice \textbf{completo} se ogni successione di Cauchy in esso definita è convergente.


Uno spazio \textbf{euclideo} è uno spazio normato la cui norma è definita attraverso un \textbf{prodotto scalare}, ovvero un'applicazione $ (\cdot , \cdot) \taleche V \times V \rightarrow \C $ che, $\forall x,y,z \in V$ e $\forall \lambda, \mu \in \C$, soddisfa
\begin{align}
	&(x,y) = \overline{(y,x)} \quad &\text{Hermitianità}\\
	&(x, \lambda y + \mu z) = \lambda(x,y) + \mu (x,z) \quad &\text{Linearità}\\
	&(x,y) \geq 0 \spacer (x,x) = 0 \leftrightarrow x=0
\end{align}

Uno spazio euclideo completo si dice \textbf{spazio di Hilbert}.

\subsection{Operatori}

Presi due spazi di Hilbert $H_1$ e $H_2$, definiamo un \textbf{operatore lineare} $A$ come un'applicazione del tipo
\begin{align}
	&A \taleche \mathcal{D}_A \subseteq H_1 \rightarrow R_A \subseteq H_2 \\
	&A(\alpha v + \beta w) = \alpha A(v) + \beta A(w) \spacer \forall v,w \in \mathcal{D}_A \quad \forall \alpha, \beta \in K
\end{align}
Un operatore si dice
\begin{enumerate}
	\item \textbf{Limitato} se $\exists N\in \R \taleche ||A(x)|| \leq N||x|| \quad \forall x \in \mathcal{D}_A$
	\item \textbf{Continuo} se $\forall \{ x_n\} \in \mathcal{D}_A \taleche x_n \rightarrow x$ segue che $\limit{n}{infty} A(x_n) = A(x)$
		\textbf{Nota bene:} un operatore lineare $(A,\mathcal{D}_A)$ è continuo se e solo se è limitato.
\end{enumerate}

\newpage

Iniziamo definendo il seguente funzionale lineare
\begin{align}
	&x \rightarrow \phi_\nu(x)\\
	&\phi_\mu(x) = (\mu, Ax) \quad \forall x \in \mathcal{D}_A
\end{align}
Dato quindi un operatore $A$, si definisce il suo \textbf{aggiunto} $A^\dagger$ come l'operatore per cui, grazie anche al teorema di Riesz
\begin{align}
	(\mu, Ax) = (A^\dagger \mu, x)
\end{align}
Se $A$ è lineare, limitato e densamente definito ne segue che $||A|| = ||A^\dagger||$.

Un'operatore $A$ si dice \textbf{unitario} quando
\begin{enumerate}
	\item è \textbf{Isometrico}, ovvero 
		\begin{align}
			&(Ax,Ay) = (x,y) \quad \forall x,y \in \mathcal{D}_A\\
			&||A|| = 1
		\end{align}
	\item il suo range è denso in H
\end{enumerate} 
Per gli operatori unitari vale che
\begin{align}
	U^\dagger U = \mathbb{1}
\end{align}

Un operatore si dice \textbf{simmetrico} quando
\begin{align}
	(x,Ay) = (Ax,y) \spacer \forall x,y \in \mathcal{D}_A
\end{align}
Se un operatore è simmetrico e $\mathcal{D}_A = \mathcal{D}_{A^\dagger}$ allora si dice che è \textbf{autoaggiunto}.


\subsubsection{Spettri per operatori Infinito-Dimensionali}

Sia $H$ uno spazio di Hilbert. Dato un \textbf{operatore lineare densamente definito}
\begin{align}
	A : \mathcal{D}_A \subseteq H \rightarrow R_A \subseteq H
\end{align} 
e definita la famiglia di operatori
\begin{align}
	&T_z(A) \taleche \mathcal{D}_A \rightarrow H\\
	&T_z(A) = z\mathbb{1} - A \spacer z \in \C
\end{align}


Si diche che $z$ appartiene all'\textbf{insieme risolvente} $Res(A)$ dell'operatore $A$ se $T_z (A)$ è \textbf{biunivoco} con \textbf{inverso limitato}.

Tale inverso prende il nome di \textbf{Operatore risolvente} di $A$ e viene definito come
\begin{align}
	&R_z(A) \taleche \mathcal{R}_{T_z(A)} \rightarrow H\\
	&R_z(A) = (z\mathbb{1} - A)^{-1}
\end{align}

Se invece $z \notin Res(A)$, allora si dice che esso appartiene allo spettro di $A$.
Rispetto agli operatori finito-dimensionali questo non implica che $z$ sia autovalore di $A$. Ci troviamo quindi a dividere gli elementi dello spettro in tre categorie:
\begin{enumerate}
	\item Lo \textbf{spettro puntuale} $\sigma_p(A)$, ovvero l'insieme degli autovalori di $A$, per i quali quindi esiste un $x \neq 0 \taleche Ax = zx$
	 
	\item Lo \textbf{spettro residuo} $\sigma_\rho(A)$, ovvero l'insieme degli $z\notin \sigma_p(A)$ per i quali $R_z(A)$ non è densamente definito.
	
	\item \textbf{Spettro continuo} $\sigma_c(A)$, ovvero l'insieme degli $z\notin \sigma_p(A) \cup \sigma_\rho (A)$  per i quali $R_z(A)$, per quanto densamente definito, non è limitato.
\end{enumerate}

\newpage

\section{Operatori integrali (Davide Bufalini, Alessandro Marcelli)}

\subsection{Spunti di teoria (Alessandro Marcelli)}

Si definiscono gli \textbf{operatori integrali} quella famiglia di operatori lineari del tipo
\begin{align}
	K \taleche L_2 [a,b] \rightarrow L_2 [a,b] 
\end{align}
per i quali l'azione su un vettore $f \in H$ è definita attraverso un \textbf{nucleo integrale} $K(x,t)$ nel seguente modo
\begin{align}
	[Kf](x) =  \int_{a}^{b} dy \; K(x,y) f(y)  
\end{align}
Definita la \textbf{norma di Hilbert-Schmidt} come
\begin{align}
	||K||_{HS} =\int_{a}^{b} dx \; \int_{a}^{b} dy \; |K(x,y)|^2
\end{align}
Gli operatori definiti attraverso nuclei per i quali tale norma è finita sono detti \textbf{nuclei di HIlbert-Schmidt} e sono \textbf{limitati}.

Le equazioni di Fredholm sono un noto caso di studio di problemi agli autovalori Per operatori integrali lineari. Si tratta di cercare soluzioni per equazioni del tipo
\begin{align}
	&\phi(x) - \lambda \int_{a}^{b} dy \; K(x,y) \phi(y) = f(x) 
\end{align}
L'equazione può essere riscritta come
\begin{align}
	&(\mathbb{1} - \lambda K)\phi=f \nextpassage
	&\left(\frac{1}{\lambda}\mathbb{1} - K \right)\phi=\frac{f}{\lambda}
\end{align}
Definendo $\mu = \lambda^{-1}$ e $T_\mu = (\mu \mathbb{1} - K)$ otteniamo 
\begin{align}
	&\left(\mu\mathbb{1} - K \right)\phi= T_\mu \phi = \tilde{f}
\end{align}
La sua omogenea associata è l'equazione agli autovalori per un operatore $K$
\begin{align}
	&\left(\mu\mathbb{1} - K \right)\phi= T_\mu \phi = 0
\end{align}
I $\lambda_i$ per cui questa equazione ammette soluzioni non triviali vengono detti \textbf{numeri caratteristici} dell'equazione, e sono l'inverso degli autovalori $\mu_i$

\newpage



\subsection{Esempio 1 (Davide Bufalini)}

Dato l'operatore integrale il cui nucleo in $L^2([0,1]x[0,1])$ è dato da
\begin{align}
	k(x,y) = 2xy-4x^2 \label{nucleo}
\end{align}
andiamo a
\begin{enumerate}
	\item discuterne lo spettro puntuale
	\item risolvere, al variare di $\lambda \in \C$ e $\alpha \in \R$, l'equazione di Fredholm di 2a specie
	\begin{align}
		\phi(x) -\lambda \int_{0}^{1} dy \; k(x,y) \phi(y) = 1 -\alpha x
	\end{align} 
\end{enumerate}

Iniziamo dal secondo punto. Possiamo riscrivere l'equazione di Fredholm come
\begin{align}
	\phi(x) &= \lambda \int_{0}^{1} dy \; (2xy-4x^2) \phi(y) + 1 -\alpha x = \nonumber\\
	&= 2\lambda x \int_{0}^{1} dy \; y\phi(y) - 4x^2 \lambda \int_{0}^{1} dy \; \phi(y) + 1 -\alpha x \label{lol}
\end{align}
Notiamo come il nucleo \ref{nucleo} sia di \textbf{Pincherle-Goursat} (o \textbf{degenere}), ossia può essere riscritto come 
\begin{align}
	k(x,y) &= \sum_{i=1}^{N} P_i(x)Q_i(y)   
\end{align}
Da cui segue che possiamo definire
\begin{align}
	A_{ij} &= \int_{0}^{1} dy \; P_i(x)Q_i(y)\\
	B_i &= \int_{0}^{1} dy \; Q_i(y)f(y) \firstpassage
	\phi(x) &= f(x)+ c_A\lambda A + c_B\lambda B \label{loll2}
\end{align}
Confrontando la \ref{loll2} e la \ref{lol} troviamo 
\begin{align}
	&A = \int_{0}^{1} dy \; \phi(y)\;\; \spacer c_A = -4x^2\\
	&B = \int_{0}^{1} dy \; y\phi(y) \spacer c_B = +2x \firstpassage
	&\phi(x) =1 -\alpha x -4x^2 \lambda A + 2x\lambda B 
\end{align}
Sostituendo in queste espressioni la \ref{lol} otteniamo
\begin{align}
	A &= \int_{0}^{1} dy \; 2y\lambda B - 4y^2 \lambda A + 1 -\alpha y = \lambda B - \frac{4}{3} A \lambda +1 - \frac{\alpha}{2}\\
	B &= \int_{0}^{1} dy \; 2y^2\lambda B - 4y^3 \lambda A + y -\alpha y^2 = \frac{2}{3}\lambda B - A \lambda +\frac{1}{2} - \frac{\alpha}{3}
\end{align}
Queste equazioni formano il seguente sistema
\begin{align}
	\double{A\left(1+\frac{4}{3} \lambda\right) -\lambda B &= 1 -\frac{\alpha}{2}}{\lambda A + \left(1 - \frac{2}{3} \lambda \right)B &= \frac{1}{2} - \frac{\alpha}{3}}
\end{align}
La cui matrice è dunque
\begin{align}
	M = \begin{pmatrix}
		1+\frac{4}{3} \lambda && - \lambda \\
		\lambda && 1-\frac{2}{3} \lambda
	\end{pmatrix}
\end{align}

Inziamo studiando l'omogenea, ovvero il caso $\det M=0$
\begin{align}
	&\det M = 0 \firstpassage
	&\left( 1+\frac{4}{3} \lambda \right) \left( 1-\frac{2}{3} \lambda \right) + \lambda^2 =0 \nextpassage
	& \frac{1}{9} \lambda^2 + \frac{2}{3} \lambda +1 = 0\nextpassage
	& \left( 1+\frac{1}{3} \lambda \right)^2 = 0 \nextpassage
	& \lambda_c = -3 \nextpassage
	&\sigma_p(k) = \left\{ -\frac{1}{3} \right\}
\end{align}

Passiamo ora al caso non omogeneo. Dobbiamo distinguere due casi
\begin{enumerate}
	\item  $\lambda \neq \lambda_c$
		\begin{align}
			A &= \frac{1}{\det (M)} \begin{vmatrix} 
	1 - \frac{\alpha}{2} && -\lambda \\
	\frac{1}{2} - \frac{\alpha}{3} && 1 - \frac{2}{3}\alpha		
	\end{vmatrix} = \dots = \frac{1 - \frac{\alpha}{2} - \frac{1}{6} \lambda}{\frac{\lambda^2}{9} + \frac{2}{3}\lambda + 1} \\
	B &= \frac{1}{\det (M)} \begin{vmatrix} 
	1 + \frac{4}{3}\alpha && 1 - \frac{\alpha}{2} \\
	\lambda && \frac{1}{2} - \frac{\alpha}{3}		
	\end{vmatrix} = \dots = \frac{\frac{1}{2} - \frac{\alpha}{3} + \frac{3\alpha -2}{6} \lambda}{\frac{\lambda^2}{9} + \frac{2}{3}\lambda + 1}	
	\end{align}
	Otteniamo quindi
	\begin{align}
		\phi(x) = 1 -\alpha x + 2 \left[ \frac{\frac{1}{2} - \frac{\alpha}{3} + \frac{3\alpha -2}{6} \lambda}{\frac{\lambda^2}{9} + \frac{2}{3}\lambda + 1} \right] \lambda x - 4 \left[ \frac{1 - \frac{\alpha}{2} - \frac{1}{6} \lambda}{\frac{\lambda^2}{9} + \frac{2}{3}\lambda + 1} \right] \lambda x^2
	\end{align}

	\item  $\lambda = \lambda_c$
	
	Prima di procedere dobbiamo trovare $\phi_0^{(+)}$, la soluzione dell'omogenea di $k^\dagger$ che servirà quando imporremo $<\phi_0^{(+)}, \phi> = 0$ tramite la quale potremo ricavare i valori $\alpha$ tali da restituire altre soluzioni. 
	
	Andiamo quindi a risolvere l'omogenea per
	\begin{align}
		k^\dagger (x,y) = \overline{k(x,y)} \overset{\R}{=} k(y,x) = 2yx - 4y^2
	\end{align} 
	Risolvendo Fredholm otteniamo quindi
	\begin{align}
		\phi_0^{(+)}(x) &= \lambda \int_{0}^{1} dy \; (2xy-4y^2) \phi_0^{(+)} = \nonumber\\
		&=  2\lambda x\int_{0}^{1} dy \; y \phi_0^{(+)}(y) -4\lambda \int_{0}^{1} dy \;y^2 \phi_0^{(+)}(y) = \nonumber\\
		&= 2\lambda x D - 4 \lambda C \spacer \double{C &= \int_{0}^{1} dy \; y \phi_0^{(+)}(y)}{D &= \int_{0}^{1} dy \;y^2 \phi_0^{(+)}(y)} \label{zero}
	\end{align}
	la \ref{zero} ci porta a scrivere che
	\begin{align}
		C &= \int_{0}^{1} dy \; \left( y^3D - 4\lambda C y^2 \right) = \frac{1}{2} \lambda D - \frac{4}{3} \lambda C\\
		D &= \int_{0}^{1} dy \; \left( 2\lambda y^2D - 4\lambda C y \right) = \frac{2}{3} \lambda D - 2 \lambda C
	\end{align}
	Da cui otteniamo il sistema
	\begin{align}
		\double{\left(1 + \frac{4}{3} \lambda\right) C - \frac{1}{2} \lambda D &= 0}{2\lambda C + \left( 1 - \frac{2}{3}\lambda \right) D &=0} \spacer M = \begin{bmatrix}
			1 + \frac{4}{3} \lambda && - \frac{1}{2} \lambda\\
			2 \lambda && 1- \frac{2}{3} \lambda
		\end{bmatrix}
	\end{align}

	L'omogenea ha lo stesso risultato di prima, con $\lambda_c = -3$, e per tale valore il sistema diventa
	\begin{align}
		&\double{\left(1 -4 \right) C + \frac{3}{2} D &= 0}{-6 C + \left( 1 +2 \right) D &=0} \firstpassage
		 &\double{D &= 2C}{C &= ap} \spacer ap = \text{a piacimento}
	\end{align}
	Ponendo $C = \frac{1}{4\lambda_c}$ otteniamo $D = \frac{1}{2 \lambda_c}$ e arriviamo alla formula
	\begin{align}
		\phi_0^{(+)}(x) &= 2D \lambda_c x - 4 \lambda_c C = x-1
	\end{align}
	Possiamo quindi procedere
	\begin{align}
		&<\phi_0^{(+)},f> = 0 \firstpassage
		&\int_{0}^{1} dx \; (x-1) (1-\alpha x) =0\nextpassage
		&\alpha = 3
	\end{align}
	Per tale valore di $\alpha$ abbiamo quindi ulteriori soluzioni al problema di partenza. Se andiamo a sostituire $\lambda = \lambda_c = -3$ otteniamo
	\begin{align}
		&\double{\left(1 -4 \lambda\right) A +3B &= 1 - \frac{3}{2}}{-3A + (1+2)B &= \frac{1}{2}-1} \firstpassage
		&\double{-3A +3B &= - \frac{1}{2}}{-3A + 3B &= -\frac{1}{2}} \nextpassage
		&\double{B = &A - \frac{1}{6}}{A = &ap} \spacer ap= \text{a piacere} \nextpassage
		&\phi_A(x) = 1 -3x +2 (-3)\left(A - \frac{1}{6}\right)x -4Ax^2(-3) \nextpassage
		&\phi_A(x) = 1 -2x -6Ax +12Ax^2
	\end{align}
\end{enumerate}

\newpage

\subsection{Esempio 2 (Davide Bufalini)}

Determinare i numeri caratteristii dell'operatore integrale il cui nucleo in $L^2([-1,+1]x[-1,+1])$ è dato da
\begin{align}
	k(x,y) = xy -x^2y^2
\end{align}
Si calcolino inoltre, per tutti i $\lambda\in\C$ le soluzioni dell'equazione di Fredholm
\begin{align}
	&\phi(x) -\lambda \int_{-1}^{+1} dy \; k(x,y) \phi(y) = 5x^3 + x^2 - 3x\\
	&g(x) = 5x^3 + x^2 - 3x
\end{align}
Nel nostro caso avremo
\begin{align}
	\phi(x) &=\lambda \int_{-1}^{+1} dy \; (xy -x^2y^2) \phi(y)  + g(x)	= \nonumber\\
	&=\lambda x \int_{-1}^{+1} dy \; \phi(y)  - \lambda x^2\int_{-1}^{+1} dy \; y^2 \phi(y) + g(x)\\
	A &= \int_{-1}^{+1} dy \; \phi(y) = \\ 
	  &=\int_{-1}^{+1} dy \; (\lambda y^2 A -\lambda y^2 B) + \int_{-1}^{+1} dy \; (5x^3 + x^2 - 3x) = \nonumber \\ 
	  &=\frac{2}{3} A \lambda + \cancel{2} -\cancel{2}\\
	B &= \int_{-1}^{+1} dy \; y^2 \phi(y) = \dots = -\frac{2}{5} \lambda B + \frac{2}{5}    
\end{align}
Abbiamo quindi il sistema
\begin{align}
	\double{\left( 1 - \frac{2}{3}\lambda \right) A &=0}{\left( 1 + \frac{2}{5}\lambda \right) B &= \frac{2}{5}}
\end{align}

Per trovare lo spettro puntuale andiamo a risolvere il caso omogeneo
\begin{align}
	&\double{\left( 1 - \frac{2}{3}\lambda \right) A &=0}{\left( 1 + \frac{2}{5}\lambda \right) B &= 0} \to \double{\lambda_1 = +\frac{3}{2}}{\lambda_2 = -\frac{5}{2}} \firstpassage
	&\sigma_p(k) = \left\{ \frac{2}{3} , -\frac{2}{5} \right\}
\end{align}
Andiamo ora a studiare il sistema non omogeneo. Distinguiamo i due casi
\begin{enumerate}
	\item $\lambda \neq \lambda_c$
	\begin{align}
		\double{\left( 1 - \frac{2}{3}\lambda \right) A &=0}{\left( 1 + \frac{2}{5}\lambda \right) B &= 0} \to \double{A &=0}{B &= \frac{2}{5+2\lambda}}
	\end{align}
	Da cui otteniamo la soluzione
	\begin{align}
		\phi(x) = g(x) - \frac{2}{5+2\lambda} x^2
	\end{align}
	\newpage
	\item $\lambda = \lambda_c$
	\begin{enumerate}
		\item $\lambda=\frac{3}{2}$
		\begin{align}
			\double{0\cdot A &=0}{\frac{8}{5} B &= \frac{2}{5}} \to \double{A = &ap}{B = &\frac{2}{8}} \spacer ap = \text{a piacere}
		\end{align}
		Da cui otteniamo la famiglia di soluzioni
		\begin{align}
			\phi_{\frac{3}{2}}(x) = g(x) - \frac{3}{8} x^2 + \frac{3}{2}xA 
		\end{align}
		\item $\lambda=-\frac{5}{2}$
		\begin{align}
			\double{0\cdot A &=0}{0\cdot B &= \frac{2}{5}} \to \text{sistema non risolvibile}
		\end{align}
	
	\end{enumerate}
\end{enumerate}

\newpage

\section{Operatori differenziali (Davide Bufalini, Alessandro Marcelli)}

\subsection{Spunti di teoria (Alessandro Marcelli)}

Gli operatori differenziali $\mathcal{L}_x^{(n)}$di ordine $n$ sono definiti come
\begin{align}
	&\double{\mathcal{L}_x^{(n)} u(x) = f(x)}{\text{Boundary conditions}}\\
	&\mathcal{L}_x^{(n)} = \sum_{k=0}^{n} a_k(x) \frac{d^k}{dx^k}
\end{align}


I problemi di ricerca di soluzioni particolari si dividono in due famiglie
\begin{align}
	&\left\{
	\begin{array}{ccc}
		\mathcal{L}_x^{(n)} u(x) &= &f(x)\\
		\frac{d^n}{dx^n}u(x_0) &= &u^{(n)}_0\\
		\frac{d^{n-1}}{dx^{n-1}}u(x_0) &= &u^{(n-1)}_0\\
		\vdots\\
		\frac{d}{dx}u(x_0) &= &u'_0\\
		u(x_0) &= &u_0\\
	\end{array}
	\right. \quad &\text{Problemi di Cauchy}\\
	&\left\{
	\begin{array}{ccc}
		\mathcal{L}_x^{(n)} u(x) &= &f(x)\\
		u(x_n) &= &u_n\\
		u(x_{n-1}) &= &u_{n-1}\\
		\vdots\\
		u(x_1) &= &u_1\\
		u(x_0) &= &u_0\\
	\end{array}
	\right. \quad &\text{Problemi di Sturm-Liouvlille}
\end{align}

La ricerca delle soluzioni può passare attraverso il calcolo del nucleo $G(x,y)$ dell'operatore inverso, denominato \textbf{funzione di Green}, per il quale vale
\begin{align}
	\mathcal{L}_x^{(n)} G(x,y) = \delta(x-y)
\end{align}
E che quindi si può definire
\begin{align}
	G(x,y) = c_1 u_1(x) + c_2 u_2(x) + \theta(x-y) \left[ \frac{u_1(y)u_2(x) - u_1(x)u_2(y)}{W(y)} \right]
\end{align}
Dove al denominatore abbiamo il \textbf{denominatore della matrice di Wronsky}, definito come
\begin{align}
	W = \det \begin{pmatrix}
		u_1^0 && u_2^0\\
		\dot{u}_1^0 && \dot{u}_2^0
	\end{pmatrix}
\end{align}
Due famiglie importanti di funzioni di Green sono le
\begin{enumerate}
	\item \textbf{Funzioni di Green avanzate}, utili nello studio dei problemi di Cauchy
	\begin{align}
		G(x,y) \equiv 0 \quad x>y
	\end{align}
	\item \textbf{Funzioni di Green ritardate}
	\begin{align}
		G(x,y) \equiv 0 \quad x<y
	\end{align}
\end{enumerate}

\newpage


\subsection{Esempio 1: Operatore Differenziale del IIo Ordine (Davide Bufalini)}

Sia l'operatore
\begin{align}
	\mathcal{L}_x^\lambda = -\frac{d^2}{dx^2} - \lambda \spacer \lambda= \text{cost.}
\end{align}

Si richiede di
\begin{enumerate}
	\item calcolare, al variare di $\lambda$, la soluzione del problema omogeneo
	\begin{align}
		\triple{\mathcal{L}_x^\lambda f(x) &= 0}{f(0) &=0}{\dot{f}(\pi) &=0}
	\end{align}
	\item ricavare la Funzione di Green (l'operatore risolvente) per i valori in cui $\mathcal{L}_x^\lambda$ è invertibile del corrispontente problema di Sturm-Liouville
\end{enumerate}

L'operatore è in forma canonica con coefficienti costanti nella forma
\begin{align}
	a \ddot{f} + b \dot{f} + cf=0 \spacer \text{nel nostro caso } b=0
\end{align}
Prendendo un Ansatz del tipo $f(x) \sim e^{\alpha x}$ otteniamo
\begin{align}
	-\alpha^2 - \lambda =0 \to \alpha= \pm i \sqrt{\lambda}
\end{align}
Da cui segue
\begin{align}
	f(x) = A e^{i \sqrt{\lambda}} + B e^{-i \sqrt{\lambda}}
\end{align}
Imponendo l condizioni di bordo otteniamo
\begin{align}
	&\double{f(0) =0}{\dot{f}(\pi) =0} \to \double{A+B = &0}{i\sqrt{\lambda} \left(A e^{i \sqrt{\lambda}} - B e^{-i \sqrt{\lambda}}\right) = &0} \firstpassage
	&\double{B = &-A}{Ai \sqrt{\lambda} 2 \cos(\sqrt{\lambda} \pi) =&0}
\end{align}
Escludendo le soluzioni banali $A=0=B$ e $\lambda=0$ abbiamo $n\in\Z$ autovalori nella forma
\begin{align}
	\lambda = \frac{(2k+1)^2}{4}
\end{align}
Per trovare gli autovettori corrispondenti ci affidiamo alla condizione $A = -B$ dalla quale segue
\begin{align}
	f_n(x)=-2i B \sin\left(\frac{2k+1}{2} x\right)
\end{align}
Per trovare B sfruttiamo l'ortogonalità
\begin{align}
	\delta_{n,m} = <f_n,f_m> &= \int_{0}^{\pi} dx \; \overline{f_n}(x) f_m(x) = \nonumber \\
	&= \int_0^\pi dx\; (2i B^*) \sin\left(\frac{2k+1}{2} x\right) \cdot (-2i B) \sin\left(\frac{2k+1}{2} x\right)
\end{align}
Iniziamo studiando il caso $n=m$
\begin{align}
	\delta_{m,m} = <f_n,f_m> &= 4|B|^2 \int_0^\pi dx\; \sin^2\left(\frac{2k+1}{2} x\right) = 4|B|^2 \frac{\pi}{2}
\end{align}
Per ortonormalità imponiamo
\begin{align}
	4|B|^2 \frac{\pi}{2} = 1 \to |B| = \frac{1}{\sqrt{2\pi}} 
\end{align}
Che, nel caso $n\neq m$, diventa 	
\begin{align}
	B = \frac{e^{i\rho}}{\sqrt{2\pi}} \spacer \rho \in \R	
\end{align}
Da cui segue che
\begin{align}
	f_n(x) = -i \sqrt{\frac{2}{\pi}}e^{i\rho}\sin \left( \frac{2n+1}{2}x \right) \spacer \double{\rho \in \R}{n\in\Z}
\end{align}

Siccome $\lambda=0$ non è autovalore non abbiamo modi nulli.

Andiamo ora a calcolare l'\textbf{operatore risolvente}.
\begin{align}
	&(\mathcal{L}_{SL} -\lambda \mathbb{1}) G_\lambda (x,y) = \delta (x-y) \spacer \lambda \neq \lambda_n \; \text{per evitare poli semplici} \firstpassage
	&-\frac{d^2}{dx^2} G_x(x,y) - \lambda G_\lambda = \delta (x-y)
\end{align}

Utilizziamo un trucco. Siccome il problema è della tipologia
\begin{align}
	-\frac{d^2f}{dx^2} - \lambda f = 0
\end{align}
Andiamo a cercare due soluzioni $f^0_{1,2}$ tali che ciascua soddisfi una delle due condizioni al bordo. Una volta trovate andiamo a calcolare il Wronskiano, necessario per la definizione della funzione di Green
\begin{align}
	G_\lambda (x,y) = \double{\frac{f_1^0(x)f_2^0(x)}{W} \spacer x<y}{\frac{f_1^0(y)f_2^0(y)}{W} \spacer x>y}
\end{align}
In base alle condizioni di bordo proviamo con
\begin{align}
	f_1^0(x) &= \sin(\sqrt{\lambda} x) \to f_1^0(0) =0\\
	f_2^0(x) &= \cos(\sqrt{\lambda} (\pi-x)) \to f_2^0(\pi) =0
\end{align}
Andiamo quindi a calcolarne il Wronskiano
\begin{align}
	W = \begin{vmatrix}
		f_1^0 && f_2^0\\
		\dot{f_1^0} && \dot{f_2^0}
	\end{vmatrix} = \begin{vmatrix}
	\sin(\sqrt{\lambda} x) && \cos(\sqrt{\lambda} (\pi-x))\\
	\sqrt{\lambda}\cos(\sqrt{\lambda}x) && \sqrt{\lambda} \cos (\sqrt{\lambda} (\pi-x))
\end{vmatrix} = \cdots = -\sqrt{\lambda} \cos( \sqrt{\lambda} \pi)
\end{align}
Da cui otteniamo
\begin{align}
	G_\lambda (x,y) = \double{-\frac{ \sin(\sqrt{\lambda} x)\cos(\sqrt{\lambda} (\pi-x))}{\sqrt{\lambda} \cos( \sqrt{\lambda} \pi)} \spacer x<y}{-\frac{ \sin(\sqrt{\lambda} y)\cos(\sqrt{\lambda} (\pi-y))}{\sqrt{\lambda} \cos( \sqrt{\lambda} \pi)} \spacer x>y}
\end{align}

Vediamo come per i $\lambda_n$ che abbiamo trovato in precedenza ci siano solo poli semplici.

\newpage

\subsection{Esempio 2 (Davide Bufalini)}

Dato l'operatore differenziale
\begin{align}
	\mathcal{L}_x = -i \frac{d}{dx} + \frac{\beta}{x} \spacer \beta \in \R
\end{align}
Il cui nucleo è dato da
\begin{align}
	\mathcal{D}_{\mathcal{L}_x} = \left\{ f \in L^2 \left( \left[\frac{1}{2},1\right] \right) \taleche f\left(\frac{1}{2}\right) = f(1) \right\}
\end{align}

Si richiede di
\begin{enumerate}
	\item Costruire $(\mathcal{L}_x^\dagger, \mathcal{D}_{\mathcal{L}_x^\dagger})$
	\item Discutere autovalori e autovettori di $(\mathcal{L}_x, \mathcal{D}_{\mathcal{L}_x})$
	\item Trovare i valori di $\beta$ per cui $(\mathcal{L}_x, \mathcal{D}_{\mathcal{L}_x})$ è autoaggiunto e/o invertibile
\end{enumerate}

Partiamo dall'\textbf{aggiunto}:
\begin{align}
	<g,\mathcal{L}_x f> &= \int_{\frac{1}{2}}^{1} dx \; \overline{g} \left(-i \frac{d}{dx} + \frac{\beta}{x}\right) f = \nonumber\\
	&= \int_{\frac{1}{2}}^{1} dx \; \overline{g} \left(-i \frac{d}{dx}f\right) + \int_{\frac{1}{2}}^{1} dx \; \overline{g}\left(\frac{\beta}{x} f\right) = \nonumber\\
	&=-i \int_{\frac{1}{2}}^{1} dx \; \overline{g} \dot{f} + \int_{\frac{1}{2}}^{1} dx \; \overline{g}\frac{\beta}{x}f = \nonumber\\
	&= \left. -i f \overline{g} \right|_{\frac{1}{2}}^{1} +i \int_{\frac{1}{2}}^{1} dx \;  \dot{\overline{g}}f + \int_{\frac{1}{2}}^{1} dx \; \overline{\left(\frac{\beta}{x} g\right)} f
\end{align}
Imponendo che $\mathcal{L}_x^\dagger$ sia autoagiunto otteniamo che
\begin{align}
	&-i f\left( \frac{1}{2} \right) \cdot\left[ \overline{g}\left(\frac{1}{2}\right) - \overline{g}(1) \right] = 0 \firstpassage
	&\overline{g}\left(\frac{1}{2}\right) = \overline{g}(1) \nextpassage
	&\mathcal{L}_x^\dagger = -i \frac{d}{dx} + \frac{\beta}{x} \spacer \mathcal{D}_{\mathcal{L}_x^\dagger} = \left\{ g \in L^2 \left( \left[\frac{1}{2},1\right] \right) \taleche g\left(\frac{1}{2}\right) = g(1) \right\}
\end{align}
Andiamo ora a calcolare gli autovalori e gli autovettori:
\begin{align}
	&\mathcal{L}_x f= \lambda f \firstpassage
	&\left( -i \frac{d}{dx} + \frac{\beta}{x} \right) f = \lambda f\nextpassage
	&\frac{df}{dx} = i\left( \lambda - \frac{\beta}{x} \right) f \nextpassage
	&\int_{\frac{1}{2}}^{x} dt \; \frac{\dot{f}(t)}{f(t)} ) = \int_{\frac{1}{2}}^{x} dt \; \left( \lambda - \frac{\beta}{x} \right) \nextpassage
	&\vdots\nextpassage
	&f(x) = f\left(\frac{1}{2}\right) e^{i\left[ \lambda \left( x - \frac{1}{2} \right)  - \beta \ln(2x)\right]}
\end{align}
Applicando le condizioni al contorno otteniamo
\begin{align}
	&f(1) =  f\left(\frac{1}{2}\right) \firstpassage
	&f\left(\frac{1}{2}\right) e^{i\frac{\lambda}{2}  - \beta \ln(2)} = f\left(\frac{1}{2}\right) e^{i\cdot 0}\nextpassage
	&e^{i\frac{\lambda}{2}  - \beta \ln(2)} = 1 \nextpassage
	&e^{i\frac{\lambda_n}{2}  - \beta \ln(2)} = e^{i2n\pi} \nextpassage
	&\frac{\lambda_n}{2}  = \beta \ln(2) + 2n\pi \nextpassage
	&\lambda_n = 2\beta \ln(2) + 4n\pi \to f_n(x) = f\left(\frac{1}{2}\right) e^{i\left[ (2\beta \ln(2) + 4n\pi) \left( x - \frac{1}{2} \right) - \beta \ln(2x)\right]}
\end{align}

\newpage

\section{Operatori a blocchi (Davide Bufalini, Alessandro Marcelli, Paolo Proia)}
\subsection{Esempio 1 con metodo esplicito e analitico (Davide Bufalini, Paolo Proia)}
Dato il seguente operatore a blocchi
\begin{align}
	\triple{(Tx)_1 \quad\;\; &= &0 \qquad \qquad \qquad}{(Tx)_{2n} \quad &= &\alpha x_{2n+1} \quad n> 1}{(Tx)_{2n+1} &= &\alpha x_{2n} \quad\quad n> 1} \spacer \alpha \in \C
\end{align}
Si richiede di
\begin{enumerate}
	\item Costruire $T^\dagger$ e calcolare $||T||$ e  $||T^\dagger||$
	\item Discutere gli spettri puntuali e relative molteplicità
	\item Discutere gli spettri residui
	\item Discutere per quali valori di $\alpha \in \C$ (se esistono) l'operatore $T$ è una combinazione delle seguenti
	\begin{enumerate}
		\item Autoaggiunto
		\item Unitario
		\item Invertibile
	\end{enumerate}
\end{enumerate}

Un possibile modo per il calcolo dell'aggiunto è iniziare esplicitando l'azione dell'operatore
\begin{align}
	\begin{matrix}
		Tx_1 &= &0\\
		Tx_2 &= &\alpha x_3\\
		Tx_3 &= &\alpha x_2\\
		Tx_4 &= &\alpha x_5\\
		Tx_5 &= &\alpha x_4
	\end{matrix} \to T= 
	\begin{bmatrix}
		0 && 0 && 0 && 0 && 0 && 0 && \dots \\
		0 && 0 && \alpha && 0 && 0 && 0 && \dots \\
		0 && \alpha && 0 && 0 && 0 && 0 && \dots \\
		0 && 0 && 0 && 0 && \alpha && 0 && \dots \\
		0 && 0 && 0 && \alpha && 0 && 0 && \dots \\
		\vdots && \vdots && \vdots && \vdots && \vdots && \vdots && \ddots
	\end{bmatrix}
\end{align}
Da cui segue che
\begin{align}
	T^\dagger= 
	\begin{bmatrix}
		0 && 0 && 0 && 0 && 0 && 0 && \dots \\
		0 && 0 && \alpha^* && 0 && 0 && 0 && \dots \\
		0 && \alpha^* && 0 && 0 && 0 && 0 && \dots \\
		0 && 0 && 0 && 0 && \alpha^* && 0 && \dots \\
		0 && 0 && 0 && \alpha^* && 0 && 0 && \dots \\
		\vdots && \vdots && \vdots && \vdots && \vdots && \vdots && \ddots
	\end{bmatrix} \to 	\triple{(T^\dagger x)_1 \quad\;\; &= &0 \;\qquad \qquad \qquad}{(T^\dagger x)_{2n} \quad &= &\alpha^* x_{2n+1} \quad n> 1}{(T^\dagger x)_{2n+1} &= &\alpha^* x_{2n} \quad\quad n> 1}
\end{align}

Un altro modo per il calcolo dell'aggiunto passa dalla sua definizione, che ci dice che
\begin{align}
	&(y,Tx) = (T^\dagger y,x) \firstpassage
	&\sum_{i=1}^{N} y^*_i (Tx)_i = \sum_{i=1}^{N} (T^\dagger y)_i^* x_i
\end{align}
Esplicitando le somme si ottengono i seguenti confronti
\begin{align}
	&0 + y_2^* \cdot \alpha x_3 + y_3^* \cdot \alpha x_2 + y_4^* \cdot \alpha x_5 + y_5^* \cdot \alpha x_4 + \dots =  \nonumber \\
	=&(T^\dagger y)_1^* \cdot x_1 + (T^\dagger y)_2^* \cdot x_2 +(T^\dagger y)_3^* \cdot x_3 +(T^\dagger y)_4^* \cdot x_4 + (T^\dagger y)_5^* \cdot x_5 + \dots\firstpassage
	&y_1 \cdot 0 = 0  = (T^\dagger y)_1^* \cdot x_1 \to (T^\dagger y)_1 = 0\\
	&\double{y_2^* \cdot (Tx)_2 = y_2^* \cdot \alpha x_3 = (T^\dagger y)_3^* \cdot x_3}{y_3^* \cdot (Tx)_3 = y_3^* \cdot \alpha x_2 = (T^\dagger y)_2^*\cdot x_2} \to  \double{(T^\dagger y)_2^* = \alpha y_3^*}{(T^\dagger y)_3^* = \alpha y_2^*} \to  \double{(T^\dagger y)_2
		= \alpha^* y_3}{(T^\dagger y)_3 = \alpha^* y_2}\\
	&\double{y_4^* \cdot (Tx)_4 = y_4^* \cdot \alpha x_5 = (T^\dagger y)_5^*\cdot x_5}{y_5^* \cdot (Tx)_5 = y_5^* \cdot \alpha x_4 = (T^\dagger y)_4^*\cdot x_4} \to \double{(T^\dagger y)_4^* = \alpha y_5^*}{(T^\dagger y)_5^* = \alpha y_4^*}\to \double{(T^\dagger y)_4 = \alpha^* y_5}{(T^\dagger y)_5 = \alpha^* y_4}\\
	&\vdots \nonumber
\end{align}
Da cui otteniamo (per fortuna) lo stesso risultato
\begin{align}
	\triple{(T^\dagger x)_1 \quad\;\; &= &0 \;\qquad \qquad \qquad}{(T^\dagger x)_{2n} \quad &= &\alpha^* x_{2n+1} \quad n> 1}{(T^\dagger x)_{2n+1} &= &\alpha^* x_{2n} \quad\quad n> 1}
\end{align}




Andiamo quindi a calcolare la norma dell'operatore
\begin{align}
	||Tx||^2 &= <Tx,Tx> = \sum_{n=1}^{\infty} \overline{(Tx)}_n (Tx)_n = \nonumber\\
			 &=\cancel{\overline{(Tx)}_1 (Tx)_1} + \sum_{n=2}^{\infty} (|\alpha|^2 |x_{2n+1}|^2 + |\alpha|^2 |x_{2n}|^2) = \nonumber\\
			 &=\sum_{n=2}^{\infty} |\alpha|^2( |x_{2n+1}|^2 + |x_{2n}|^2) = \nonumber\\
			 &=|\alpha|^2 \sum_{n=2}^{\infty} |x_{n}|^2 =|\alpha|^2( ||x||^2 - |x_1|^2) \leq |\alpha|^2 \sum_{n=1}^{\infty} |x_{n}|^2 = |\alpha|^2 ||x||^2 \firstpassage
	||T|| &\leq |\alpha|	
\end{align}
Se $|x_1|=0$ allora la norma è raggiunta(?).
Siccome per operatori limitati e densamente definiti allora $||T|| = ||T^\dagger||$ e quindi 
\begin{align}
	||T^\dagger|| &\leq |\alpha|	
\end{align}
Andiamo ora a calcolare lo \textbf{spettro puntuale}. Dobbiamo risolvere
\begin{align}
	T\vec{x} &= \lambda \vec{x}	
\end{align}
Da cui otteniamo
\begin{align}
	\begin{matrix}
		Tx_1 =& 0 &= \lambda x_1 \\
		Tx_2 =& \alpha x_3 &= \lambda x_2\\
		Tx_3 =& \alpha x_2 &= \lambda x_3\\
		Tx_4 =& \alpha x_5 &= \lambda x_4\\
		Tx_5 =& \alpha x_4 &= \lambda x_5\\
		&\dots
	\end{matrix} 
\end{align}
\newpage
Distinguiamo ora i vari casi
\begin{enumerate}
	\item $\lambda = 0 \spacer \alpha =0$
		\begin{align}
			0 &= 0\cdot x_1 \to x_1 = \text{qualsiasi} \\
			0 &= 0\cdot x_2 \to x_2=\text{qualsiasi}\\
			0 &= 0\cdot x_3 \to x_3=\text{qualsiasi} \\
			&\dots \firstpassage
			\vec{x} &= (x_1,x_2, x_3, \dots)
		\end{align}
	\item $\lambda = 0 \spacer \alpha \neq 0$
		\begin{align}
			0 &= 0\cdot x_1 \to x_1 = \text{qualsiasi} \\
			\alpha x_3 &= 0\cdot x_2 \to x_3=0\\
			\alpha x_2 &= 0\cdot x_3 \to x_2=0 \\
			&\dots \firstpassage
			\vec{x} &= (x_1,0, 0, \dots)
		\end{align}
	\item $\lambda \neq 0 \spacer \alpha =0$
		\begin{align}
			0 &= \lambda x_1 \to x_1=0 \\
			0 &= \lambda x_2 \to x_3=0\\
			0 &= \lambda x_3 \to x_2=0 \\
			&\dots \firstpassage
			\vec{x} &= (0,0, 0, \dots)
		\end{align}
	\item $\lambda \neq 0 \spacer \alpha \neq 0$
		\begin{align}
			0 &= \lambda x_1 \to x_1=0 \\
			\alpha x_2 &= \lambda x_3 \to x_2=\frac{\lambda}{\alpha} x_3\\
			\alpha x_3 &= \lambda x_2 \to x_3=\frac{\lambda^2}{\alpha^2} x_3 \to \left(1 - \frac{\lambda^2}{\alpha^2}\right) x_3 = 0 \\
			\alpha x_4 &= \lambda x_5 \to x_4=\frac{\lambda}{\alpha} x_5\\
			\alpha x_5 &= \lambda x_4 \to x_5=\frac{\lambda^2}{\alpha^2} x_5 \to \left(1 - \frac{\lambda^2}{\alpha^2}\right) x_5 = 0\\
			&\dots \firstpassage
			x_{2n} &=\frac{\lambda}{\alpha} x_{2n+1} \to \left(1 - \frac{\lambda^2}{\alpha^2}\right) x_{2n+1} = 0 
		\end{align}
	
		\newpage
		Dobbiamo distinguere due casi
		\begin{enumerate}
			\item $|\lambda| \neq |\alpha|$
				\begin{align}
					&\left(1 - \frac{\lambda}{\alpha^2}\right) x_{2n+1} = 0 \to x_{2n+1} = 0 \;\; \forall n \in \N \to x_n = 0 \forall n\in \N \firstpassage
					&\vec{x} = \vec{0}
				\end{align}
			\item $|\lambda| = |\alpha| \to \lambda = e^{i\phi} |\alpha|$
				\begin{align}
					&x_{2n} =\frac{e^{i\phi} |\alpha|}{|\alpha|} x_{2n+1} \to x_{2n} = e^{i\phi}x_{2n+1}\firstpassage
					&\vec{x} = (0, e^{i\phi}x_{3}, x_3, e^{i\phi}x_{5}, x_5,\dots)\\
					&\sigma_p(T) = \{ \lambda \in \C \quad|\quad \lambda = e^{i\phi} |\alpha| \spacer \alpha \in \C \}
				\end{align}
		\end{enumerate}
\end{enumerate}

Si ottiene un risultato analogo per $\sigma(T^\dagger)$.

Andiamo ora a calcolare lo \textbf{spettro residuo} dell'operatore e del suo aggiunto.

Andiamo quindi a cercare 
\begin{align}
	&z\neq \lambda \taleche \exists \vec{\eta} \neq \vec{0} \taleche <\vec{\eta} , \vec{v}> =0 \\
	&\vec{v} = (z\mathbb{1} - T)\vec{x} \in \text{Range}(z\mathbb{1} - T) \quad \forall \vec{x}\in \mathcal{D}_T
\end{align}
Nel nostro caso abbiamo
\begin{align}
z\mathbb{1} - T= 
	\begin{bmatrix}
		z      && 0       && 0       && 0       && 0       && 0      && \dots\\
		0      && z       && -\alpha && 0       && 0       && 0      && \dots\\
		0 	   && -\alpha && z       && 0       && 0       && 0      && \dots\\
		0      && 0       && 0       && z       && -\alpha && 0      && \dots\\
		0 	   && 0       && 0       && -\alpha && z       && 0      && \dots\\
		\vdots && \vdots  && \vdots  && \vdots  && \vdots  && \vdots && \ddots
	\end{bmatrix}
\end{align}
Da cui segue che
\begin{align}
	(z\mathbb{1} - T)\vec{x} = \left( \begin{matrix}
		zx_1\\
		zx_2 - \alpha x_3\\
		zx_3 - \alpha x_2\\
		zx_4 - \alpha x_5\\
		zx_5 - \alpha x_4\\
		\vdots
	\end{matrix} \right)
\end{align}
E arriviamo quindi a scrivere
\begin{align}
	0 &= <\vec{\eta}, \vec{v} > = \sum_{n} \overline{\eta}_n v_n = \nonumber\\
	&= \eta_1^*zx_1 + \eta_2^*(zx_2 - \alpha x_3) + \eta_3^*(zx_3 - \alpha x_2) + \dots = \nonumber \\
	&= x_1z\eta_1^* + x_2(z\eta_2^* - \alpha \eta_3^*) + x_3(z\eta_3^* - \alpha \eta_2^*) + \dots
\end{align}
Siccome la somma è nulla, questo ci porta a scrivere
\begin{align}
	\eta^*_1 =& 0 \\
	\eta^*_2 =& \frac{\alpha}{\lambda} \eta^*_3\\
	\eta^*_3 =& \frac{\alpha}{\lambda} \eta^*_2 = \frac{\alpha^2}{\lambda^2} \eta^*_3 \to  \left(1-\frac{\alpha^2}{\lambda^2}\right) \eta^*_3 = 0\\
	&\vdots \nonumber
\end{align}
Abbiamo però un "problema". Siccome $z \neq \lambda = e^{i\phi} |\alpha|$, abbiamo che
\begin{align}
	1-\frac{\alpha^2}{\lambda^2} \neq 0 \to \eta_3^* = 0 \to \vec{\eta} = \vec{0}
\end{align}
Ma questo non è possibile per definizione, quindi abbiamo
\begin{align}
	\sigma_\rho(T) = \emptyset
\end{align}

In conclusione rispondiamo alle domande del quarto punto:
\begin{enumerate}
	\item L'operatore è \textbf{autoaggiunto} per $\alpha = \alpha^*$, ovvero per $\alpha \in \R$, dato che in tale caso segue $T= T^\dagger$ e $\mathcal{D}_T=\mathcal{D}_{T\dagger}$
	\item Non può mai essere \textbf{unitario}, dato che $TT^\dagger \neq \mathbb{1} \neq T^\dagger T$
	\item Non è \textbf{invertibile} dato che $z\in \sigma_p(T)$
\end{enumerate}


\newpage

\subsection{Esempio 2 con metodo esplicito e analitico (Davide Bufalini, Alessandro Marcelli)}

Sia l'operatore in $l^2(\C)$ definito da
\begin{align}
	\double{(Tx)_1 =& 2^a x_2\quad\quad\quad\quad\quad\quad\quad\quad\;\;\;\;\;}{(Tx)_n =& x_n + (n+1)^ax_{n+1} \quad n \geq 2}
\end{align}

Si richiede di
\begin{enumerate}
	\item Costruire l'aggiunto
	\item Discuttere gli spettripuntuali
	\item Discutere se, al variare di $a\in\R$, $z_0=0$ appartiene allo spettro dell'operatore o dell'aggiunto
\end{enumerate}


Iniziamo dal primo punto.

Un possiible metodo è quello di \textbf{esplicitare} l'azione dell'operatore
\begin{align}
	\begin{matrix}
		Tx_1 &= &2^a x_2\\
		Tx_2 &= &x_2 +3^a x_3\\
		Tx_3 &= &x_3 +4^a x_4\\
		\vdots
	\end{matrix} \to T= 
	\begin{bmatrix}
		0 	   && 2^a    && 0 	   && 0 	 && 0 	   && 0 	 && \dots \\
		0 	   && 1 	 && 3^a    && 0 	 && 0 	   && 0 	 && \dots \\
		0	   && 0      && 1 	   && 4^a 	 && 0 	   && 0 	 && \dots \\
		0 	   && 0 	 && 0	   && 1	     && 5^a    && 0	     && \dots \\
		0 	   && 0 	 && 0 	   && 0		 && 1 	   && 6^a 	 && \dots \\
		\vdots && \vdots && \vdots && \vdots && \vdots && \vdots && \ddots
	\end{bmatrix}
\end{align}
Da cui otteniamo quindi
\begin{align}
T^\dagger= 
	\begin{bmatrix}
		0 	   && 0      && 0 	   && 0 	 && 0 	   && 0 	 && \dots \\
		2^a    && 1 	 && 0      && 0 	 && 0 	   && 0 	 && \dots \\
		0	   && 3^a    && 1 	   && 0 	 && 0 	   && 0 	 && \dots \\
		0 	   && 0 	 && 4^a	   && 1	     && 0      && 0	     && \dots \\
		0 	   && 0 	 && 0 	   && 5^a	 && 1 	   && 0 	 && \dots \\
		\vdots && \vdots && \vdots && \vdots && \vdots && \vdots && \ddots
	\end{bmatrix} \to T^\dagger = \double{(T^\dagger x)_1 =& 0}{(T^\dagger x)_n =& x_n + n^a x_{n-1} \quad n \geq 2}
\end{align}
Altrimenti possiamo calcolarlo \textbf{analiticamente}, studiando il prodotto
\begin{align}
	&(y,Tx) = (T^\dagger y, x)\firstpassage
	&\sum_{i=1}^{n} y_i^* (Tx)_i =  \sum_{i=1}^{n} (T^\dagger y)_i^* x_i
\end{align}
Da cui otteniamo che
\begin{align}
	y_1^* (Tx)_1 &= y_1^* 2^a x_2 = (T^\dagger y)_1^* x_1\\
	y_2^* (Tx)_2 &=	y_2^*( x_2 + 3^a x_3) = (T^\dagger y)_2^* x_2\\
	y_3^* (Tx)_3 &=	y_3^*( x_3 + 4^a x_4) = (T^\dagger y)_3^* x_3\\
	y_4^* (Tx)_4 &=	y_4^*( x_4 + 4^a x_5) = (T^\dagger y)_4^* x_4\\
	y_5^* (Tx)_5 &=	y_5^*( x_5 + 6^a x_6) = (T^\dagger y)_5^* x_5\\
	&\vdots \nonumber
\end{align}
Possiamo quindi scrivere
\begin{align}
	&y_1^* 2^a x_2 + y_2^* x_2 + y_2^* 3^a x_3 + y_3^* x_3 + y_3^* 4^a x_4  + y_4^* x_4 + y_4^* 5^a x_5 + \dots = \nonumber \\
	=& (y_1^* 2^a + y_2^*)x_2 + (y_2^* 3^a + y_3^*) x_3 + (y_3^* 4^a + y_4^*) x_4 =  \\
	=&(T^\dagger y)_1^* x_1 + (T^\dagger y)_2^* x_2 + (T^\dagger y)_3^* x_3 + (T^\dagger y)_4^* x_4 + \dots
\end{align}
Confrontando i termini otteniamo quindi
\begin{align}
	(T^\dagger y)_1^* &= 0\\
	(T^\dagger y)_2^* &= y_1^* 2^a + y_2^*\\
	(T^\dagger y)_3^* &= y_2^* 3^a + y_3^*\\
	(T^\dagger y)_4^* &=y_3^* 4^a + y_4^*\\
	\vdots \nonumber
\end{align}
Da cui (per fortuna) otteniamo di nuovo
\begin{align}
	 T^\dagger = \double{(T^\dagger x)_1 =& 0}{(T^\dagger x)_n =& x_n + n^a x_{n-1} \quad n \geq 2}
\end{align}

Andiamo ora a calcolare gli spettri puntuali dell'operatore e del suo aggiunto.
\begin{align}
	&Tx = \lambda x \to (Tx)_n = \lambda x_n\firstpassage
	&(Tx)_1 = 2^a x_2 = \lambda x_1\\
	&(Tx)_2 = x_2 + 3^a x_3 = \lambda x_3\\
	&(Tx)_3 = x_3 + 4^a x_4 = \lambda x_4\\
	&(Tx)_4 = x_4 + 5^a x_5 = \lambda x_5\\
	\vdots \nonumber
\end{align}

Dobbiamo ora distinguere i diversi casi: 

\begin{enumerate}

	\item $\lambda = 0 \spacer a=0$
		\begin{align}
			&x_1 = \text{indeterminato}\\
			&x_2 = 0\\
			&x_2 + x_3 = 0 \to x_3 =0\\
			&x_3 + x_4 = 0 \to x_4 =0\\
			&x_4 + x_5 = 0 \to x_5=0\\
			&\vdots \nextpassage
			&\vec{x} = (x_1, 0, 0, 0, \dots) \spacer \nu = 1
		\end{align}
		
	\item $\lambda = 0 \spacer a \neq 0$
		\begin{align}
			&x_1 = \text{indeterminato}\\
			&2^a x_2 = 0 \to x_2 = 0\\
			&x_2 + 3^a x_3 = 0 \to x_3 =0\\
			&x_3 + 4^a x_4 = 0 \to x_4 =0\\
			&x_4 + 5^a x_5 = 0 \to x_5=0\\
			&\vdots \nextpassage
			&\vec{x} = (x_1, 0, 0, 0, \dots) \spacer \nu = 1
		\end{align}

	\item $\lambda \neq 0 \spacer a \neq 0$
		\begin{align}
			&2^ax_2 = \lambda x_1 \to \double{x_1 =& \text{indeterminato}}{x_2 =& \frac{\lambda}{2^a} x_1}\\
			&x_2 + 3^a x_3 = \lambda x_2 \to x_3 = \frac{\lambda -1}{3^a} x_2 \\
			&x_3 + 4^a x_4 = \lambda x_3 \to x_4 = \frac{\lambda -1}{4^a} x_3 = \frac{(\lambda -1)^2}{(3\cdot4)^a} x_2  = \frac{\lambda(\lambda -1)^2 }{(2\cdot 3\cdot4)^a} x_1 =  \frac{\lambda(\lambda -1)^2 }{(4!)^a} x_1\\
			&x_4 + 5^a x_5 = \lambda x_4 \to x_5 = \frac{\lambda -1}{5^a} x_4 = \dots =  \frac{\lambda(\lambda -1)^3 }{(5!)^a} x_1\\
			&\vdots \nextpassage
			&\triple{x_1 =& \text{indeterminato}}{x_2 =& \frac{\lambda}{2^a} x_1}{x_n =& \frac{\lambda(\lambda -1)^{n-2} }{(n!)^a} x_1 \spacer n>2}
		\end{align}

		La situazione si complica. Dobbiamo trovare per quali $\lambda$ e $a$ abbiamo $\vec{x}\in l^2(\C)$. Andiamo a calcolare la norma
		\begin{align}
			||x||^2 &= \sum_{n=1}^{+\infty} |x_n|^2 = |x_1|^2 + \left| \frac{\lambda}{2^a} x_1 \right|^2 + \sum_{n=3}^{+\infty} \left| \frac{\lambda(\lambda -1)^{n-2} }{(n!)^a} x_1 \right|^2 = \nonumber \\
			&= \left( 1 + \frac{|\lambda|^2}{2^{2a}} \right) |x_1|^2 + \sum_{n=3}^{+\infty}  \frac{|\lambda|^2|\lambda -1|^{2n-4} }{(n!)^{2a}} \left| x_1 \right|^2
		\end{align}
		Applichiamo ora il criterio della radice, ovvero
		\begin{align}
			n| \sim \frac{n^n}{e^n} \sqrt{2\pi n}
		\end{align}
		E otteniamo
		\begin{align}
			\frac{|\lambda|^2|\lambda -1|^{2n-4} }{(n!)^{2a}} \sim \frac{|\lambda|^\frac{2}{n}|\lambda -1|^\frac{{2n-4}}{n} }{\left(\frac{n^n}{e^n} \sqrt{2\pi n}\right)^\frac{2a}{n}} \sim \left| \frac{(\lambda -1)^2 e^{2a}}{n^{2a}} \right| \to \triple{0 \quad &a>0}{|\lambda -1|^2 \quad &a=0}{\infty \quad &a<0}
		\end{align}
		Il caso $a<0$ diverge, e quindi lo escludiamo a priori. Studiamo ora gli altri due
		\begin{enumerate}
			\item $a = 0 \to |\lambda -1|^2<1 \to |\lambda -1| < 1$ 
			
			Abbiamo quindi una regione di convergenza circolare centrata in 1 con raggio di convergenza pari a 1.
			
			\item $a > 0 \to 0$
			
			In questo caso la serie converge $\forall \lambda \in \C$, e abbiamo un raggio di convergenza infinito.
		\end{enumerate}
		
\end{enumerate}

\newpage

Andiamo ora a studiare lo spettro puntuale dell'aggiunto.
\begin{align}
	&T^\dagger v = \lambda v \to (T^\dagger  v)_n = \lambda v_n\firstpassage
	&(T^\dagger  v)_1 = 0 = \lambda v_1\\
	&(T^\dagger  v)_2 = v_2 + 2^a v_1 = \lambda v_2\\
	&(T^\dagger  v)_3 = v_3 + 3^a v_2 = \lambda v_3\\
	&(T^\dagger  v)_4 = v_4 + 4^a v_3 = \lambda v_4\\
	\vdots \nonumber
\end{align}

Di nuovo, andiamo a distinguere i vari casi
\begin{enumerate}
	\item $\lambda = 0 \spacer a=0$
		\begin{align}
			&v_1 = \text{qualsiasi}\\
			&v_2 + v_1 = 0 \to v_2 = -v_1\\
			&v_3 + v_2 = 0 \to v_3 = -v_2 = v_1\\
			&v_4 + v_3 = 0 \to v_4 = -v_3 = -v_1\\
			\vdots \nextpassage
			&\vec{v} = (v_1, -v_1 , v_1, \dots, (-1)^{n+1}v_1, \dots) \spacer \nu = 1
		\end{align}
		

	
	\item $\lambda = 0 \spacer a \neq 0$
		\begin{align}
			&v_1 = \text{qualsiasi}\\
			&v_2 + 2^a v_1 = 0 \to v_2 = -2^a v_1 \\
			&v_3 + 3^a v_2 = 0 \to v_3 = -3^a v_2 = 6^a v_1 \\
			&v_4 + 4^a v_3 = 0 \to v_2 = -4^a v_3 = -24^a v_1  \\
			\vdots \nextpassage
			&\vec{v} = (v_1, -2^a v_1 ,6^a v_1, \dots, (-1)^{n+1} (n!)^a v_1, \dots) \spacer \nu = 1
		\end{align}
		


	\item $\lambda \neq 0 \spacer a \neq 0$
		\begin{align}
			&v_1 = 0\\
			&v_2 = \lambda v_2 \to (1 - \lambda)v_2 = 0 \to \double{v_2 = 0}{\lambda = 1}\\
			&(1 - \lambda)v_3 = -3^a v_2\\ 
			\vdots \nonumber
		\end{align}
		Distinguiamo i due casi
		\begin{enumerate}
			\item $\lambda \neq 1 \to v_2 = 0 \to v_3 = 0 \to \dots \to \vec{v} = \vec{0}$
			\item $\lambda = 1 \to v_2 = \text{qualsiasi} \to v_3 = v_3 -3^a v_2 \to v_2 = 0 \to \vec{v} = \vec{0}$
		\end{enumerate}
\end{enumerate}

\newpage

Andiamo ora a studiare lo \textbf{spettro residuo}. Vale il teorema per cui
\begin{align}
	&\sigma_\rho (A) \subset \overline{\sigma_\rho (A^\dagger)} \leftrightarrow \sigma_\rho (A^\dagger) \subset \overline{\sigma_\rho (A)}\\
	&0 \in\sigma_\rho (A) \to 0 \in \overline{\sigma_\rho (A^\dagger)}
\end{align}

CHIEDERE A PAOLO

\chapter{Equazioni Differenziali Ordinarie}



\section{Spunti di teoria generale (Alessandro Marcelli)}

\subsection{Concetti generali}

Un'\textbf{equazione differenziale ordinaria} (EDO) di ordine $n$ viene genricamente definita come
\begin{align}
	\Phi(x, u, u', u'', \dots, u^{(n)}) = 0
\end{align}
In questo corso ci si occupa solo di equazioni che possono essere scritte in \textbf{forma normale}, ovvero dove si può esplicitare in funzione della derivata massima
\begin{align}
	u^{(n)}=f(x, u, u', u'', \dots, u^{(n-1)})
\end{align}
con la proprietà di poter essere riscritte come sistemi di $n$ EDO del 1o ordine
nel seguente modo
\begin{align}
	\left\{\begin{array}{c}
		v_1 \quad = u' \hfill \\
		v_2 \quad = v_1' = u'' \hfill \\
		\dots \hfill \\
		v_{n-1} = v_{n-2}' \hfill \\
		v_n \quad = f(x, u, v_1, v_2, \dots, v_{n-1}) \hfill 
	\end{array}
	\right.
\end{align}
Questo sistema definisce un insieme infinito di curve detto \textbf{integrale generale} dell'equazione.

Per individuare una curva specifica bisogna specificare almeno un punto  per cui deve passare, e i rispettivi valori che le varie derivate vi devono assumere fino alla $u^{(n-1)}$. Questo prende il nome di \textbf{problema di Cauchy}.

Nota: visto che come abbiamo visto i problemi di ordine superiore al primo possono essere visti come combinazione di EDO al primo ordine, nella teoria ci appoggeremo a problemi del 1o ordine del tipo
\begin{align}
	\left\{\begin{array}{c}
		u'(x) \; = F(x,u(x)) \hfill \\
		u(x_0) \;  = u_0 \hfill 
	\end{array}
	\right.
\end{align}

\subsection{Esistenza ed unicità delle soluzioni}

L'esistenza delle soluzioni di una EDO è dettata dai seguenti due teoremi:
\begin{enumerate}
	\item \textbf{Teorema di Cauchy:} \textit{sia assegnato un problema di Cauchy in un dominio
		\begin{align}
			\mathcal{D} = (a,b)x(c,d) \taleche (x_0,u_0) \in \mathcal{D}
		\end{align}
		Se rispetto ad $u$ si verifica la condizione di Lipschitzianità per $F(x,u)$
		\begin{align}
			|F(x,u_1) - F(x,u_2)| \leq L|u_1 - u_2|
		\end{align}
		Allora la soluzione del problema di Cauchy esiste ed è unica.	
	}
	\item \textbf{Teorema di Peano:} \textit{sia assegnato un problema di Cauchy in un dominio
		\begin{align}
			\mathcal{D} = (a,b)x(c,d) \taleche (x_0,u_0) \in \mathcal{D}
		\end{align}
		Se la $F(x,u)$ è continua, allora esiste un opportuno intorno $(x_0, x_0+\delta) \in \mathcal{D}$ tale per cui la soluzione del problema esiste.
	}
\end{enumerate}

Qualora ci si trovi nel secondo caso, non è più sufficiente un unico punto per determinare la soluzione, e bisogna ijnvece specificare, in numero pari all'ordine della EDO, delle \textbf{condizioni di bordo} (in inglese \textbf{Boundary Conditions}, BC), arriviamo così a definire due tipi di problemi:
\begin{enumerate}
	\item \textbf{Problemi di Dirichlet:} condizioni al bordo su $u(x)$. 
	
	Nel caso di equazioni lineari, questi si riconducono allo studio degli autovalori dei corrispondenti operatori lineari, e prendono il nome di \textbf{Problemi di Sturm-Liouville}.
	\item \textbf{Problemi di Neumann:} condizioni al bordo su $u'(x)$
\end{enumerate}

\subsection{Distribuzioni}

Le \textbf{distribuzioni} sono una famiglia di operatori lineari integrali definite come
\begin{align}
	(T,\phi) &= \int_{-\infty}^{+\infty} dx \; T(x) \phi(x)\\
	(T',\phi) &= \int_{-\infty}^{+\infty} dx \; T'(x) \phi(x) = -\int_{-\infty}^{+\infty} dx \; T(x) \phi'(x) = -(T,\phi')
\end{align}

Per quanto riguarda questo corso ci limiteremo allo spazio $S^{\infty}$ delle \textbf{funzioni di prova}, con derivate continue di ogni ordine che all'infinito tendono a zero più rapidamente di ogni potenza.

Esempi noti di distribuzioni sono la \textbf{delta} e la \textbf{theta}.

La continuità delle distribuzioni come funzionali consente di ottenere una loro definizione come limite di altre distribuzioni.

\subsection{Equazioni Differenziali Lineari}
Si parla di EDO lineari di ordine $n$ quando possono essere poste nella forma
\begin{align}
	&\mathcal{L}^{(n)}_x u(x) = f(x) \label{edolin}\\
	&\mathcal{L}^{(n)}_x = \sum_{k=0}^{n} c_k(x) \frac{d^k}{dx^k}
\end{align}
Lo studio passa quindi da $F$ ai vari $c_k$. Notiamo come possiamo avere singolarità della soluzione solo intorno a punti singolari dei coefficienti, che definiamo come \textbf{punti singolari fissi}.

Spesso le soluzioni sono note solo attraverso sviluppi in serie di Laurent intorno a tali punti.

Nel caso lineare, la continuità dei coefficienti ci garantisce esistenza ed unicità delle soluzioni.

Definiamo la \textbf{omogenea associata} dell'equazione \ref{edolin} come
\begin{align}
	&\mathcal{L}^{(n)}_x u(x) = \sum_{k=0}^{n} c_k(x) \frac{d^k}{dx^k}u(x) = 0
\end{align}
L'integrale generale di una EDO lineare può essere scritto come combinazione della soluzione dell'omogenea $u_{om}(x)$ e di una soluzione particolare $u_{p}(x)$
\begin{align}
	u(x) = u_{om}(x) + u_{p}(x)
\end{align}


\section{EDO lineari del 1o ordine [teoria](Alessandro Marcelli)}

Sia il problema di Cauchy
\begin{align}
	\left\{\begin{array}{c}
		c_1 (x) u'(x) + c_0(x) u(x) \; = f(x) \hfill \\
		u(x_0) \;  = u_0 \hfill 
	\end{array}
	\right. \label{ode1ord1}
\end{align}
Possiamo risolvere il problema in due modi:

\subsection{Separazione di variabili + variazione delle costanti}
	\begin{enumerate}
		\item La soluzione dell'omogenea si ricava per \textbf{separazione delle variabili}
		\begin{align}
			&c_1 (x) u'(x) + c_0(x) u(x) \; = 0 \firstpassage
			&c_1 (x) u'(x) = - c_0(x) u(x) \nextpassage
			&\frac{du}{dx} = -\frac{c_0(x)}{c_1(x)} u(x)\nextpassage
			&\int_{u_0}^u du \; \frac{1}{u}= - \int_{x_0}^x dt \;  \frac{c_0(t)}{c_1(t)}\nextpassage
			&\ln(u(x)) - C = - \int_{x_0}^x dt \;  \frac{c_0(t)}{c_1(t)} \spacer C = \ln (u_0)\nextpassage
			&u_{om}(x) = e^C e^{- \int_{x_0}^x dt \;  \frac{c_0(t)}{c_1(t)}}
		\end{align}
		\item La soluzione della particolare invece col metodo della \textbf{variazione delle costanti}, ovvero ponendo
		\begin{align}
			&u_p(x) = a(x) u_{om}(x)\firstpassage
			&u'_p(x)= a'(x) u_{om}(x) +  a(x) u'_{om}(x) \nextpassage
			&c_1(x) \cdot [a'(x) u_{om}(x) +  a(x) u'_{om}(x)] + c_0 a(x) u_{om}(x) =f(x) \nextpassage
			&c_1(x)u_{om}(x) \cdot a'(x) + \cancel{[c_1(x)u'_{om}(x) + c_0(x)u_{om}(x)]\cdot a(x)} = f(x)\nextpassage
			&c_1(x)u_{om}(x) \cdot a'(x) = f(x) \nextpassage
			&a'(x) =\frac{f(x)}{c_1(x)u_{om}(x)} \nextpassage
			&a(x) = \int_{x_0}^{x} dt\; \frac{f(t)}{c_1(t)u_{om}(t)} \nextpassage
			&u_p(x) = \left[ \int_{x_0}^{x} dt\; \frac{f(t)}{c_1(t)u_{om}(t)} \right] \cdot \left[ e^C e^{- \int_{x_0}^x dt \;  \frac{c_0(t)}{c_1(t)}} \right]
		\end{align}
		\item si sommano le due soluzioni, e si trova la curva che rispetta le condizioni iniziali del problema
	\end{enumerate}
	
\subsection{Metodo della funzione di Green} 
Il metodo della \textbf{funzione di Green} lo abbiamo già introdotto nel capitolo degli operatori. Andiamo a cercare una soluzione nel senso delle distribuzioni per
	\begin{align}
		\mathcal{L}^{(1)}_x G(x,y) = \delta(x-y) \label{green1}
	\end{align}
Abbiamo visto come l'integrale generale, la soluzione fondamentale è definita a meno di soluzioni dell'omogenea. 

Prendiamo in considerazione un problema a \textbf{condizioni omogenee} (o \textbf{nulle})
\begin{align}
	\left\{\begin{array}{c}
		c_1 (x) u'(x) + c_0(x) u(x) \; = g(x) \hfill \\
		v(x_0) \;  = 0 \hfill 
	\end{array}
	\right.
\end{align}

\textbf{Nota:} qualunque problema del tipo \ref{ode1ord1} può essere riscritto in questo modo col cambio di variabili
\begin{align}
	v(x) &= u(x) - u_0\\
	g(x) &= f(x) - c_0(x)u_0
\end{align}

Per il \textbf{teorema dell'alternativa}, se l'omogenea ammette solo la soluzione nulla, allora il problema ha soluzione unica.

In questo caso $\mathcal{L}^{(1)}_x$ risulta essere invertibile e la soluzione fondamentale con condizione iniziale omogenea $G(x_0,y) = 0$ della \ref{green1} prende il nome di \textbf{funzione di Green} dell'operatore ed è il nucleo integrale dell'operatore inverso. Grazie ad essa possiamo ricavare la soluzione come \textbf{convoluzione} col termine forzante $g(x)$
\begin{align}
	u(x) = \int dy \; G(x,y) g(y) = G * g
\end{align}
Andiamo a risolvere la \ref{green1}
\begin{align}
	&c_1 (x) \frac{d}{dx}G(x,y) + c_0(x) G(x,y)  = \delta(x-y) \firstpassage
	&\frac{d}{dx}G(x,y) + \frac{c_0(x)}{c_1(x)} G(x,y)  = \frac{\delta(x-y)}{c_1(x)}\firstpassage
	&\int_{y-\epsilon}^{y+\epsilon} dx \;\frac{d}{dx}G(x,y) + \int_{y-\epsilon}^{y+\epsilon} dx \;\frac{c_0(x)}{c_1(x)} G(x,y)  = \int_{y-\epsilon}^{y+\epsilon} dx \;\frac{\delta(x-y)}{c_1(x)}
\end{align}
Essendo i $c_k(x)$ regolari, l'integrando può al più fare un "salto" e quindi la sua primitiva $F$ è continua, e quindi nel limite $\epsilon \to 0$ il secondo integrale si annulla
\begin{align}
\limit{\epsilon}{0}\int_{y-\epsilon}^{y+\epsilon} dx \;\frac{c_0(x)}{c_1(x)} G(x,y)  = \int_{y}^{y} dx \;\frac{c_0(x)}{c_1(x)} G(x,y) = F(y,y) - F(y,y) = 0
\end{align}
Otteniamo quindi la condizione di salto
\begin{align}
&G(y+\epsilon,y) - G(y-\epsilon,y) = \frac{1}{c_1(y)}\firstpassage
&G(y_+,y) = G(y_-,y) + \frac{1}{c_1(y)}
\end{align}

\newpage
Andiamo ora a calcolare una soluzione dell'omogenea per \ref{green1}
\begin{align}
	&\frac{d}{dx}G(x,y) + \frac{c_0(x)}{c_1(x)} G(x,y)  = 0\firstpassage
	&G(x,y) = A_-(y) exp\left( - \int_{x}^{y} d\xi\; \frac{c_0(\xi)}{c_1(\xi)}\right) \; ; \;  x<y\\
	&G(x,y) = A_+(y) exp\left( - \int_{x}^{y} d\xi\; \frac{c_0(\xi)}{c_1(\xi)}\right) \; ; \; x>y
\end{align}
Applicando la condizione di salto
\begin{align}
	&A_+(y_) = A_-(y) + \frac{1}{c_1(y)} \firstpassage
	&G(x,y) = A_-(y) \cdot exp\left( - \int_{x}^{y} d\xi\; \frac{c_0(\xi)}{c_1(\xi)}\right) \; ; \;  x<y\\
	&G(x,y) = \left[  A_-(y) + \frac{1}{c_1(y)} \right] \cdot exp\left( - \int_{x}^{y} d\xi\; \frac{c_0(\xi)}{c_1(\xi)}\right) \; ; \;  x>y \firstpassage
	&G(x,y) = \left[  A_-(y) + \frac{\theta(x-y)}{c_1(y)} \right] \cdot exp\left( - \int_{x}^{y} d\xi\; \frac{c_0(\xi)}{c_1(\xi)}\right)	
\end{align}
Nel caso di problemi alle condizioni omogenee, possiamo fare l'ipotesi $x_0 < y$ , e la funzione di Green del problema deve soddisfare $G(x_0,y) = 0$, il che ci restituisce
\begin{align}
&G(x,y) = \frac{\theta(x-y)}{c_1(y)} \cdot exp\left( - \int_{y}^{x} d\xi\; \frac{c_0(\xi)}{c_1(\xi)}\right)	
\end{align}
Invece, nel caso di problemi generali del tipo \ref{ode1ord1}, il procedimento si snoda in due passaggi:
\begin{enumerate}
	\item Calcoliamo la convoluzione della generica soluzione fondamentale, tenendo a mente che $a<x_0<b$
	\begin{align}
		u(x) &= \int_a^b dy \; G(x,y) f(y) =	\nonumber\\
			 &= \int_{x_0}^b dy \;\left[  A(y) + \frac{\theta(x-y)}{c_1(y)} \right] \cdot exp\left( - \int_{y}^{x} d\xi\; \frac{c_0(\xi)}{c_1(\xi)}\right) \cdot f(y) =	\nonumber\\
			 &= \int_{x_0}^b dy \; A(y) f(y) \cdot exp\left( - \int_{y}^{x} d\xi\; \frac{c_0(\xi)}{c_1(\xi)}\right) + \nonumber\\
			 &+ \int_{x_0}^x dy \;  \frac{f(y)}{c_1(x)} \cdot exp\left( - \int_{y}^{x} d\xi\; \frac{c_0(\xi)}{c_1(\xi)}\right) 
	\end{align}
	\item Imponiamo ora le condizioni iniziali del problema, ottenendo così la definizione implicita di $A(y)$
	\begin{align}
		u_0 &= \int_{x_0}^b dy \; A(y) f(y) \cdot exp\left( - \int_{y}^{_0} d\xi\; \frac{c_0(\xi)}{c_1(\xi)}\right) 
	\end{align}
	Da cui otteniamo infine l'integrale generale
	\begin{align}
		u(x) &= u_0 \cdot  exp\left( - \int_{y}^{x} d\xi\; \frac{c_0(\xi)}{c_1(\xi)}\right) + \int_{x_0}^x dy \;  \frac{f(y)}{c_1(x)} \cdot exp\left( - \int_{y}^{x} d\xi\; \frac{c_0(\xi)}{c_1(\xi)}\right) 
	\end{align}
\end{enumerate}

\newpage

\section{EDO lineari del 2o ordine (Adriano Chialastri, Alessandro Marcelli)}


\subsection{Spunti di teoria (Alessandro Marcelli)}

Le EDO del 2o ordine sono rappresentabili come
\begin{align}
	c_2(x) u''(x) + c_1(x) u'(x) + c_0(x)u(x) = f(x)
\end{align}
Seguendo il procedimento illustrato in \ref{canon}, possono essere sempre ridotte in \textbf{forma canonica}
\begin{align}
	v''(x) + c(x)v(x) = g(x)
\end{align}
dove, ovviamente le condizioni al contorno vanno trasformate di conseguenza.

Rispetto alle EDO del 1o ordine, c'è la difficoltà aggiuntiva data dal fatto che non sempre si può trovare una soluzione dell'omogenea (e quindi l'integrale generale). Questo problema non si pone però per due famiglie di equazioni omogenee:
\begin{enumerate}
	\item \textbf{Omogenee a coefficienti costanti:}
	\begin{align}
		a u''(x) + b u'(x) + c u(x) = 0 \label{edo2ord1}
	\end{align}
	In questo caso si procede con l'\textbf{Ansatz} (o \textbf{ipotesi}) $u(x) = e^{\alpha x}$ e si va a cercare il parametro $\alpha$ tramite l'equazione caratteristica
	\begin{align}
		&a \alpha^2 + b \alpha + c = 0\firstpassage
		& \alpha_{1,2} = \frac{-b \pm \sqrt{b^2 - 4ac}}{2a} \firstpassage
		&u_1(x) = e^{\alpha_1 x} \spacer u_2(x) = e^{\alpha_2 x}
	\end{align}
	Se le soluzioni coincidono, e quindi si ha
	\begin{align}
		\alpha = -\frac{b}{2a}
	\end{align}
	si procede nel seguente modo
	\begin{align}
		&u_1(x) = e^{\alpha x} \spacer u_2(x) = \rho(x)e^{\alpha x} \firstpassage
		&\cancel{(a \alpha^2 + b \alpha + c)\rho(x)e^{\alpha x}} + (a \rho''(x) + b \rho'(x) + 2a\alpha \rho'(x))e^{\alpha x} = 0
	\end{align}
	Il primo termine si annulla, vista la \ref{edo2ord1}, otteniamo quindi
	\begin{align}
		&(a \rho''(x) + b \rho'(x) + 2a\alpha \rho'(x))e^{\alpha x} = 0\firstpassage
		&a \rho''(x) + b \rho'(x) + 2a\alpha \rho'(x)=0 \nextpassage
		&a \rho''(x) + b \rho'(x) - \cancel{2a}\frac{b}{\cancel{2a}} \rho'(x)=0 \to a \rho''(x) + \cancel{b \rho'(x)} - \cancel{b\rho'(x)} =0 \nextpassage
		&\rho''(x) = 0 \nextpassage
		&\rho(x) = c_1x + c_2
	\end{align}
	la soluzione più semplice si ha per $c_1 = 1$ e $c_2 = 0$, che ci restituisce
	\begin{align}
		u_2(x) = x e^{\alpha x}
	\end{align}

	\item \textbf{Equazioni di Eulero del 2o ordine:}
	\begin{align}
		u''(x) + \frac{a}{x}u'(x) + \frac{b}{x^2}u(x) = 0 \spacer a,b \;\;\ \text{costanti}
	\end{align}
	Stavolta le soluzioni le andiamo a cercare con l'Ansatz $u(x) = x^\alpha$, ottenendo la seguente equazione caratteristica
	\begin{align}
		&\alpha (\alpha - 1) + a\alpha + b = 0 \firstpassage
		&\alpha^2 + (a-1)\alpha + b = 0 \nextpassage
		&\alpha_{1,2} = \frac{(1-a) \pm \sqrt{(a-1)^2 - 4b}}{2} \firstpassage
		&u_1 (x) = x^{\alpha_1} \spacer u_2 (x) = x^{\alpha_2}
	\end{align}
	Nel caso di soluzioni coincidenti
	\begin{align}
		a = \frac{1-a}{2}
	\end{align}
	conviene trasformare l'equazione in una a coefficienti costanti col cambio di variabili
	\begin{align}
		&x(t) = e^t\firstpassage
		&\frac{d}{dx}u(x(t)) =e^{-t} \frac{du}{dt} \nextpassage
		&\frac{d^2}{dx^2}u(x(t)) =e^{-2t} \frac{d^2u}{dt^2} - e^{-2t} \frac{du}{dt} \nextpassage
		&e^{-2t} \left[ \frac{d^2u}{dt^2} + (a-1)\frac{du}{dt} + bu \right] = 0 \nextpassage
		&u_1(t) = e^{\alpha t} \spacer u_2(t) = te^{\alpha t} \nextpassage
		&u_1(x) = x^\alpha \spacer u_2(x) = x^\alpha \ln(x)
	\end{align}
\end{enumerate}

\newpage

\subsubsection{Metodo del Wronskiano}

Si definisce il \textbf{determinante di Wronsky} ( anche detto \textbf{Wronskiano}) di due funzioni $u_1$ e $u_2$ come
\begin{align}
	W(u_1(x),u_2(x)) = W(x) = \begin{vmatrix}
		u_1(x) & u_2(x) \\
		\dot{u}_1(x) &\dot{u}_2(x)
	\end{vmatrix} = u_1(x) \dot{u}_2(x) - \dot{u}_1(x) u_2(x) 
\end{align}
Se  $u_1(x)= a\cdot u_2(x)$, ovvero sono linearmente dipendenti, allora $W=0$.

Presa in considerazione una EDO del tipo
\begin{align}
	&a(x) \ddot{u}(x) + b(x) \dot{u}(x) + c(x) u(x) = 0 \firstpassage
	&\ddot{u}(x) + \frac{b(x)}{a(x)} \dot{u}(x) + \frac{c(x)}{a(x)} u(x) = 0\nextpassage
	&\ddot{u}(x) + a_1(x) \dot{u}(x) + a_2(x) u(x) = 0 \\
	&a_1(x) =  \frac{b(x)}{a(x)} \spacer a_2(x)=\frac{c(x)}{a(x)}
\end{align}
e prese duesaoluzioni $u_1(x)$ e $u_2(x)$ possiamo ricavare
\begin{align}
	&u_1(x)\ddot{u}_2(x) + a_1(x) u_1(x)\dot{u}_2(x) + a_2(x) u_1(x)u_2(x) = 0 \\
	&u_2(x)\ddot{u}_1(x) + a_1(x) u_2(x)\dot{u}_1(x) + a_2(x) u_2(x)u_1(x) = 0\firstpassagecomm{sottraggo membro a membro}
	&u_1(x)\ddot{u}_2(x) - u_2(x)\ddot{u}_1(x) + a_1(x) u_1(x)\dot{u}_2(x) - a_1(x) u_2(x)\dot{u}_1(x) + \cancel{ a_2(x) u_1(x)u_2(x)} - \cancel{a_2(x) u_2(x)u_1(x)} = 0 \nextpassage
	&[u_1(x)\ddot{u}_2(x) - u_2(x)\ddot{u}_1(x)]+ a_1(x) [  u_1(x)\dot{u}_2(x) - u_2(x)\dot{u}_1(x)] = 0\nextpassage
	&\dot{W}(x) + a_1(x)W(x) = 0
\end{align}
Per separazione delle variabili, e definendo $W(x_0)=W_0$ otteniamo la \textbf{formula di Liouville}
\begin{align}
	W(x) = W_0\cdot exp\left( -\int_{x_0}^{x} d\xi a_1(\xi) \right)
\end{align}
Da cui notiamo che
\begin{enumerate}
	\item In formacanonica $W=0$
	\item Se $\exists x_1 \taleche W(x_1) = 0 \to W(x) = 0\quad\forall x$
	\item Se $u_1$ e $u_2$ sono lin ind allora $W(x) \neq 0 \quad \forall x$
\end{enumerate}
Nota una delle due soluzioni, dalla formula di Liouville possiamo ricavare una definizione implicita dell'altra nel seguente modo
\begin{align}
	&\frac{W}{u_1^2(x)} = \frac{u_1(x)\dot{u}_2(x) - u_2(x)\dot{u}_1(x)}{u_1^2(x)}\\
	&\frac{W}{u_1^2(x)} = \frac{W_0\cdot exp\left( -\int_{x_0}^{x} d\xi a_1(\xi) \right)}{u_1^2(x)} \firstpassage
	&\frac{u_1(x)\dot{u}_2(x) - u_2(x)\dot{u}_1(x)}{u_1^2(x)} = \frac{W_0\cdot exp\left( -\int_{x_0}^{x} d\xi a_1(\xi) \right)}{u_1^2(x)}
\end{align}
Il termnine a sinistra altro non è che 
\begin{align}
&\frac{u_1(x)\dot{u}_2(x) - u_2(x)\dot{u}_1(x)}{u_1^2(x)} = \frac{d}{dx} \left( \frac{u_2(x)}{u_1(x)} \right)
\end{align}
Da cui otteniamo, sostituendo $W_0$ con $A$ per liberare l'estremo inferior edi integrazione
\begin{align}
&\frac{d}{dx} \left( \frac{u_2(x)}{u_1(x)} \right) = \frac{W_0\cdot exp\left( -\int_{x_0}^{x} d\xi a_1(\xi) \right)}{u_1^2(x)} \firstpassage
& \int^x d\eta \; \frac{d}{d\eta} \left( \frac{u_2(\eta)}{u_1(\eta)} \right) =\int^x d\eta \; \frac{A \cdot exp\left( -\int_{x_0}^{\eta} d\xi a_1(\xi) \right)}{u_1^2(\eta)} + c_1 \nextpassage
&u_2(x) = Au_1(x) \int^x d\eta \; \frac{1}{u_1^2(\eta)}\cdot exp\left( -\int_{x_0}^{\eta} d\xi a_1(\xi) \right) + c_1u_1(x)
\end{align}


\subsubsection{Metodo della funzione di Green}
Consideriamo ora l'operatore differenziale in forma normale
\begin{align}
	\mathcal{L}_x^{(2)} =\frac{d^2}{dx^2} + a_1(x)\frac{d}{dx} + a_2(x)
\end{align}
Ne definiamo una \textbf{soluzione fondamentale} $G(x,y)$ come
\begin{align}
	&\mathcal{L}_x^{(2)} G(x,y)= \delta(x-y) \label{greenedo2}
\end{align}
Imponiamo la condzione di continuità 
\begin{align}
	&\frac{d^2}{dx^2}G(x,y) \propto \delta(x-y)
\end{align}
Da cui 
\begin{align}
	&\left.\frac{d}{dx}G(x,y) \right|_{y-\epsilon}^{y+\epsilon} + 0 = 1 \firstpassage
	&G'(y_+,y) - G'(y_-,y) = 1
\end{align}
Siano ora $u_1(x)$ e $u_2(x)$ due soluzioni della omogenea associata alla \ref{greenedo2}, possiamo scrivere
\begin{align}
	G(x,y) = A u_1(x) + B u_2(x) + G_P(x,y)
\end{align}
La soluzone particolare può essere scritta nella forma
\begin{align}
	G_P(x,y) = \double{c_1 u_1 (x) + c_2 u_2 (x) \quad x<y}{d_1 u_1 (x) + d_2 u_2 (x) \quad x>y} \spacer c_i,d_i = \double{costanti}{c_i(y), d_i(y)}
\end{align}
Per la condizione di continuità abbiamo che
\begin{align}
	&c_1 u_1 (y) + c_2 u_2 (y) = d_1 u_1 (y) + d_2 u_2 (y)\firstpassage
	&(c_1 - d_1) u_1 (y) +(c_2 - d_2)u_2 (y) =0
\end{align}
Mentre per il salto della derivata
\begin{align}
	&d_1 \dot{u}_1 (y) + d_2 \dot{u}_2 (y) - c_1 \dot{u}_1 (y) - c_2 \dot{u}_2 (y) = 1 \firstpassage
	&(d_1 - c_1)\dot{u}_1 (y) + (d_2 - c_2)\dot{u}_2 (y) = 1 
\end{align}
Otteniamo quindi il sistema
\begin{align}
	\double{(c_1 - d_1) u_1 (y) +(c_2 - d_2)u_2 (y) =0}{(d_1 - c_1)\dot{u}_1 (y) + (d_2 - c_2)\dot{u}_2 (y) = 1}
\end{align}
Se prendiamo come incognite le differenze tra i coefficienti tra parentesi, possiamo riscrivere il sistema come il prdotto matrice-vettore
\begin{align}
	\begin{pmatrix}
		u_1 (y) & u_2 (y)\\
		\dot{u}_1 (y) & \dot{u}_2 (y)
	\end{pmatrix} \left( \begin{array}{c}
	c_1 - d_1\\
	c_2 - d_2
\end{array} \right) = \left( \begin{array}{c}
0\\
1
\end{array} \right)
\end{align}
Notiamo come la matrice dei coefficienti altro non è che la matrice di Wronsky di $u_1(x)$ e $u_2(x)$, e che quindi non ha mai determinante nullo. 

Possiamo quindi scrivere le soluzioni come
\begin{align}
d_1 - c_1 &= \frac{1}{W(y)} \cdot \begin{vmatrix}
			0 & u_2(y) \\
			1 & \dot{u}_2(y)
	\end{vmatrix} = -\frac{u_2(y)}{W(y)} \to d_1 = c_1  -\frac{u_2(y)}{W(y)}\\
d_2 - c_2 &= \frac{1}{W(y)} \cdot \begin{vmatrix}
	u_1(y) & 0\\
	\dot{u}_1(y) & 1
\end{vmatrix} = +\frac{u_1(y)}{W(y)} \to d_2 = c_2 +\frac{u_1(y)}{W(y)}
\end{align}
Trovate queste possiamo riscrivere la soluzione particolare come
\begin{align}
	G_P(x,y) = \double{c_1 u_1 (x) + c_2 u_2 (x) \hfill x<y}{\left[  c_1  -\frac{u_2(y)}{W(y)} \right] u_1 (x) + \left[ c_2 +\frac{u_1(y)}{W(y)} \right] u_2 (x) \quad x>y} \spacer c_i,d_i = \double{costanti}{c_i(y), d_i(y)}
\end{align}
Possiamo riscrivere la seconda come
\begin{align}
&\left[  c_1  -\frac{u_2(y)}{W(y)} \right] u_1 (x) + \left[ c_2 +\frac{u_1(y)}{W(y)} \right] u_2 (x)\firstpassage
&c_1 u_1 (x) + c_2 u_2 (x) + \frac{u_1(y)u_2(x) - u_1(x)u_2(y)}{W(y)}
\end{align}
Conforntando con il caso $x<y$ e impiegando il funzionale a gradino $\theta(x-y)$ possiamo riscrivere la soluzuione particolare nella seguente forma compatta
\begin{align}
	G_P(x,y)= c_1 u_1 (x) + c_2 u_2 (x) + \theta(x-y) \frac{u_1(y)u_2(x) - u_1(x)u_2(y)}{W(y)}
\end{align}
Riportiamo due casi noti che ci saranno utili più avanti
\begin{enumerate}
	\item \textbf{Funzione di Green avanzata}
		\begin{align}
			&G(x,y) \equiv 0 \quad x<y 
		\end{align}
	\item \textbf{Funzione di Green ritardata}
		\begin{align}
			&G(x,y) \equiv 0 \quad x>y 
		\end{align}
\end{enumerate}
\newpage

\subsubsection{Teorema di Green}

Il teorema di Green consente di stabilire le condizioni al bordo per il calcolo della funzione di Green di un problema differenziale lineare.

Iniziamo considerando i problemi di Sturm-Liouville con condizioni accessorie date sul bordo dell'intervallo $[a,b]$
\begin{align}
	\triple{\mathcal{L}_x^{(2)} v(x) &= g(x) \hfill}{\mathcal{\tilde{B}}_a(v) &= \tilde{\alpha}_1 v(a) + \tilde{\beta}_1 \dot{v}(a) - \gamma_1 = 0}{\mathcal{\tilde{B}}_b(v) &= \tilde{\alpha}_2 v(b) + \tilde{\beta}_2 \dot{v}(b) - \gamma_2 = 0}
\end{align}
Sappiamo dalla teoria delle EDO al 1o ordine che 
\begin{enumerate}
	\item è sempre possibile ricondurre un problema in forma canonica
	\begin{align}
		\mathcal{L}_x^{C} = \frac{d^2}{dx^2} + c(x) \label{omcond3}
	\end{align} 
	\item con l'opportuno cambio di variabili ci si può ricondurre al caso di condizioni omogenee
		\begin{align}
			\double{v(x) = u(x)+ mx + q}{\dot{v}(x) = \dot{u}(x) +m \hfill}			
		\end{align}
\end{enumerate}
Possiamo quindi riscrivere
\begin{align}
	\triple{\mathcal{L}_x^{C} u(x) &= f(x) \hfill}{\mathcal{B}^{om}_a(u) &= \alpha_1 u(a) + \beta_1 \dot{u}(a) = 0}{\mathcal{B}^{om}_b(u) &= \alpha_2 u(b) + \beta_2 \dot{u}(b) = 0} \label{omcond1}
\end{align}
ed enunciare il \textbf{teorema di Green:} \textit{Sia il problema di Sturm-Liouville appena definito. Se l'omogenea ammette solo la soluzione banale, ovvero
\begin{align}
	u^{om}(x) = 0
\end{align}
all'ora l'integrale generale è dato da
\begin{align}
	u(y) = \int_{a}^{b} dx \; G(x,y) f(x) + \left.\left[ u(x) \frac{dG(x,y)}{dx} - G(x,y) \frac{du(x)}{dx} \right]\right|_a^b
\end{align}
Dove si ha che $G(x,y)$ è soluzione fondamentale di
\begin{align}
	\mathcal{L}_x^{C} G(x,y) = \delta(x-y) \label{omcond2}
\end{align}
}
La dimostrazione passa per la funzione integrale
\begin{align}
	E(y) = \int_{a}^{b} dx \; [G(x,y) \mathcal{L}_x^{C} u(x) - u(x)  \mathcal{L}_x^{C} G(x,y)]
\end{align}
Utilizzando la \ref{omcond1} e la \ref{omcond2} otteniamo
\begin{align}
	E(y) &= \int_{a}^{b} dx \; [G(x,y) f(x) - u(x)  \delta(x-y)] = \nonumber \\
		 &= -u(y) +\int_{a}^{b} dx \; G(x,y) f(y) \label{greendim1}
\end{align}
Usando invece la \ref{omcond3} otteniamo
\begin{align}
	E(y) &= \int_{a}^{b} dx \; G(x,y) [\ddot{u}(x) + c(x)u(x)] - u(x)  [\ddot{G}(x,y) + c(x) G(x,y)] = \nonumber\\
	&= \left.[G(x,y)\dot{u}(x) - u(x)\dot{G}(x,y)]\right|_a^b - \int_{a}^{b} dx \; \cancel{\dot{G}(x,y) \dot{u}(x)} - \cancel{\dot{u}(x) \dot{G}(x,y)}  \label{greendim2}
\end{align}
Uguagliando \ref{greendim1} e \ref{greendim2} segue la tesi.

\subsubsection{Applicazione del teorema di Green ai problemi di Sturm-Liouville}

Notiamo come nel caso di condizioni
\begin{align}
	\double{u(a)=0}{u(b)=0}
\end{align}
si abbia
\begin{align}
	\double{\dot{G}(a,y)=0}{\dot{G}(b,y)=0}
\end{align}
Quindi, siccome non sono note le derivate $\dot{a}$ e $\dot{b}$, basta imporre le condizioni
\begin{align}
	\double{G(a,y)=0}{G(b,y)=0}
\end{align}
Questo può essere generalizzato al problema generale. Magari scoccio Paolo e lo vedo con lui.

Possiamo quindi dire che dato il problema
\begin{align}
	\triple{\mathcal{L}_x^{C} u(x) &= f(x) \hfill}{\mathcal{B}^{om}_a(u) &= \alpha_1 u(a) + \beta_1 \dot{u}(a) = 0}{\mathcal{B}^{om}_b(u) &= \alpha_2 u(b) + \beta_2 \dot{u}(b) = 0} \label{omcond1}
\end{align}
La funzione di Green associata sarà data da
\begin{align}
	\triple{\mathcal{L}_x^{C} G(x,y) &= \delta(x-y)}{\mathcal{B}^{om}_a(G) &= 0 \hfill}{\mathcal{B}^{om}_b(G) &= 0 \hfill}
\end{align}
Da qui possiamo ricavare una definizione operativa. Definiamo due soluzioni $u_1(x)$ e $u_2(x)$ e imponiamo
\begin{align}
	\double{\mathcal{B}^{om}_a(u_1(x)) = 0}{\mathcal{B}^{om}_b(u_2(x)) = 0}
\end{align}
Preso il Wronskiano costante $W(u_1(x),u_2(x))$ possiamo scrivere la funzione di Green come
\begin{align}
	G(x,y) = \double{\frac{1}{W(u_1(x),u_2(x))}u_1(x)u_2(y) \quad x<y}{\frac{1}{W(u_1(x),u_2(x))}u_1(y)u_2(x) \quad x>y}
\end{align}
Questa soluzione rispetta 
\begin{enumerate}
	\item le condizioni al bordo (per costruzione) 
	\item la condizione di continuità di $G(x,y)$
	\item la condizione di salto pari a 1 della derivata
\end{enumerate}
 \newpage
\subsubsection{Applicazione del teorema di Green ai problemi di Cauchy}

Ora prendiamo in considerazione le seguenti condizioni iniziali
\begin{align}
	\double{\dot{u}(x_0)=0}{u(x_0)=0}
\end{align}
Applichiamo il teorema di Green, integrando da $x_0$ a $x = \Delta$ e ponendo $\Delta > y$. 
Allora la formula data dal teorema di Green diventa
\begin{align}
	u(y) = \int_{x_0}^{\Delta} dx \; G(x,y) f(x) + \left.\left[ u(x) \frac{dG(x,y)}{dx} - G(x,y) \frac{du(x)}{dx} \right]\right|_{x_0}^\Delta
\end{align}
Scrivendo per esteso i termini al bordo dati dal secondo termine avremo
\begin{align}
	u(\Delta) \dot{G}(\Delta,y) - u(x_0) \dot{G}(x_0,y) - G(\Delta,y) \dot{u}(\Delta) + G(x_0,y) \dot{u}(x_0) 
\end{align}
Per annullare i termini in $\Delta$, che non conosciamo, dobbiamo imporre 
\begin{align}
	&\dot{G}(\Delta, y) = G(\Delta, y) = 0	
\end{align}
Ma, siccome $\Delta>y$, questa condizione altro non è che la \textbf{funzione di Green avanzata}. Si dice quindi che la funzione di Green "propaga" la soluzione da $x_0$ a $y$, e possiamo riscrivere
\begin{align}
	u(y) = \int_{x_0}^{\Delta} dx \; G(x,y) f(x) +  G(x_0,y) \dot{u}(x_0) - u(x_0) \dot{G}(x_0,y) 
\end{align}
Nel caso di condizioni omogenee, sappiamo che $G(x_0,y) = 0 = \dot{G}(x_0,y)$ e quindi ci riduciamo al problema di convoluzione
\begin{align}
	u(y) = \int_{x_0}^{\Delta} dx \; G(x,y) f(x)
\end{align}
Se non ci si vuole ricondurre al caso di condizioni omogenee, bisogna tener da conto dei termini aggiuntivi. Ma perché dovremmo volerci far del male?


\newpage

\subsection{Messaggio dall'autore}

Riporto qui, sotto richiesta dell'autore della maggior parte di questi esercizi, un suo messaggio:

\textit{Io sottoscritto Adriano Chialastri non mi riterrò responsabile di qualsiasi danno a cose o persone causato da questi esercizi, in quanto non svolti nel pieno esercizio delle mie facoltà mentali a causa nella prolungata permanenza presso il luogo correntemente noto come Sogene.
}

\subsection{Esempio 1: Problema di Cauchy (Adriano Chialastri) \label{canon}}
Risolvere il seguente problema di Cauchy
\begin{align}
	\triple{x^2 \ddot{y} - 2x \dot{y} + 2y = -x}{\dot{y}(1)=0}{y(1)=0}
\end{align}
Iniziamo passando in \textbf{forma canonica}
\begin{align}
	&y(x) = A(x)v(x)	\firstpassage
	&x^2 (\ddot{A}(x)v(x) + 2\dot{A}(x)\dot{v}(x) + A \ddot{v}(x)) - 2x (\dot{A}(x)v(x) + A \dot{v}(x)) + 2A(x)v(x) = -x \nextpassage
	&Ax^2 \ddot{v}(x) + 2(\dot{A}(x) x^2 - A(x)x) \dot{v}(x) + (\ddot{A}(x) x^2 - 2x\dot{A}(x) +2A)v(x) = -x	
\end{align}
Per poter fare il passaggio dobbiamo imporre che
\begin{align}
	&\dot{A}(x) x^2 - A(x)x = 0 \firstpassage
	&\dot{A}(x) x = A(x) \nextpassage
	&A(x) = x  
\end{align}
Da cui otteniamo che
\begin{align}
	&x^3 \ddot{v}(x) + (-\cancel{2x} +\cancel{2x})v(x) = -x	\firstpassage
	&\ddot{v}(x) = -\frac{1}{x^2} \spacer y(x) = x v(x)
\end{align}
Il nostro problema di Cauchy in forma canonica è quindi
\begin{align}
	\triple{\ddot{v}(x) = -\frac{1}{x^2}}{\dot{v}(1)=0}{v(1)=0}
\end{align}
Andiamo ora a risolvere l'\textbf{omogenea}
\begin{align}
	&\ddot{v}(x) = 0 \firstpassage
	&v_{om}(x) = Ax + B  \label{kekwdiff1}\firstpassage
	&\double{\dot{v}_{om}(1)=A=0}{v_{om}(1)=A+B=0} \to \double{A=0}{B=0} \to \text{l'operatore differenziale è invertibile} 
\end{align}
Dalla \ref{kekwdiff1} vediamo come
\begin{align}
	&v_{om}(x) = A v_1(x) + B v_2(x) \nextpassage
	&v_{om_1}(x) = x \spacer v_{om_2}(x) = 1
\end{align}
Calcoliamo il Wronskiano dell'equazione
\begin{align}
	W(v_1,v_2) = \begin{vmatrix}
		x && 1 \\
		1 && 0
	\end{vmatrix} = -1
\end{align}
E andiamo quindi a calcolare la \textbf{Funzione di Green}
\begin{align}
	G(x,y) = - \delta (x-y) (y-x) = (x-y) \delta(x-y)
\end{align}
Da cui otteniamo
\begin{align}
	v(x) &= \int_{1}^x dy \;G(x,y) \left(-\frac{1}{y^2}\right) = \nonumber \\
		 &= \int_{1}^x dy \; \delta (x-y) (y-x) \cdot \frac{1}{y^2}= \nonumber \\
		 &= \int_{1}^x dy \; \frac{y-x}{y^2}= \nonumber \\
		 &= \int_{1}^x dy \; \left(\frac{1}{y} - \frac{x}{y^2}\right)= \nonumber \\
		 &= \left[ \ln(y) + \frac{x}{y} \right]_1^x = \ln(x) +1 -x \nextpassage
	y(x) &= xv(x) = x\ln(x) - x^2 +x
\end{align}
Verifichiamo la validità della soluzione
\begin{align}
	&\dot{y}(x) = 2 - 2x + \ln(x)\\
	&\ddot{y}(x) = -2 + \frac{1}{x} \firstpassage
	&\double{\dot{y}(1) = 0}{y(1) = 0} \to \double{0 -\cancel{1} +\cancel{1} =0}{\cancel{2}-\cancel{2} + 0 =0}\\
	&x^2 \ddot{y} - 2x \dot{y} + 2y = -x \firstpassage
	&x^2 \left( -2 + \frac{1}{x} \right) -2x \left( 2 - 2x + \ln(x) \right) + 2 \left( x\ln(x) - x^2 +x \right) = -x \nextpassage
	& -\cancel{2x^2} + x - 4x + \cancel{4x^2} - 2x\ln(x) + 2x\ln(x) -\cancel{2x^2} +2x = -x \nextpassage
	& + x - 4x - \cancel{2x\ln(x)} + \cancel{2x\ln(x)} +2x = -x \nextpassage
	& -x = -x
\end{align}
E quindi l'esercizio è verificato.

\newpage

\subsection{Esempio 2: Omogenea di Eulero (Adriano Chialastri)}
Riprendiamo il precedente problema di Cauchy
\begin{align}
	\triple{x^2 \ddot{y} - 2x \dot{y} + 2y = -x}{\dot{y}(1)=0}{y(1)=0}
\end{align}
Partiamo dall'\textbf{omogenea}
\begin{align}
	&x^2 \ddot{y} - 2x\dot{y} + 2y = 0 \firstpassage
	&\ddot{y} - \frac{2}{x}\dot{y} + \frac{2}{x^2}y = 0
\end{align}
L'omogena scritta in questa forma risulta essre un'\textbf{equazione di Eulero}. Per risolverla poniamo quindi l'\textbf{Ansatz} $y=x^\alpha$ dove otteniamo
\begin{align}
	&\alpha(\alpha -1) x^{\alpha-2} - 2 \alpha x^{\alpha-2} + 2x^{\alpha-2} = 0 \firstpassage
	&\alpha^2 - \alpha - 2\alpha +2 =0 \nextpassage
	&\alpha^2 - 3\alpha +2 =0 \nextpassage
	&\alpha_{1,2} = \frac{3 \pm \sqrt{9-8}}{2} = \frac{3\pm 1}{2} = \double{2}{1} \nextpassage
	&y_{{om}_1}(x) = x \spacer y_{{om}_2}(x) = x^2 \nextpassage
	&y_{om}(x) = Ax + Bx^2   
\end{align}
Andiamo ora a cercare una \textbf{soluzione particolare} col \textbf{metodo delle costanti}. Poniamo
\begin{align}
	&y_p(x) = A(x)x + B(x)x^2 \firstpassage
	&\dot{y_p}(x)  = \dot{A}(x)x+ A(x) + \dot{B}(x) x^2 + 2B(x)x\\
	&\ddot{y_p}(x) = \ddot{A}(x)x+ 2\dot{A}(x) + \ddot{B}(x) x^2 + 4\dot{B}(x)x + 2B  \firstpassage
	&x^2 (\ddot{A}(x)x+ 2\dot{A}(x) + \ddot{B}(x) x^2 + 4\dot{B}(x)x + 2B) -2 (\dot{A}(x)x+ A(x) + \dot{B}(x) x^2 + 2B(x)x) + \nonumber \\
	&+2(A(x)x + B(x)x^2) = -x  \nextpassage
	&\dots \nextpassage
	&x^3 \ddot{A}(x) + x^4 \ddot{B}(x) + 2x^3 \dot{B}(x) = -x
\end{align}
Da cui otteniamo
\begin{align}
	&y(x) = y_p(x) + y_{om}(x) = y_p(x) + Ax + Bx^2 \firstpassage
	&\double{\dot{y}(1) = \dot{y}_p(1) + A + 2B = 0}{y(1) = y_p(1) + A + B=0}
\end{align}

\newpage

\subsection{Esempio 3: Problema di Sturm-Liouville (Adriano Chialastri)}

Risolvere il seguente problema di Sturm-Liouville
\begin{align}
	\triple{x \ddot{u}(x) + 2 \dot{u}(x) + 4xu(x) = 4}{u \left( \frac{\pi}{4}\right)=0}{u \left( \frac{\pi}{2}\right)=0}
\end{align}
Iniziamo passando il problema in \textbf{forma canonica}
\begin{align}
	&u(x) = A(x)v(x) \firstpassage
	&x(\ddot{A}(x)v(x) + 2\dot{A}(x) \dot{v}(x) + A(x) \ddot{v}(x)) + 2(\dot{A}(x)v(x) + A(x) \dot{v}(x)) + 4xA(x)v(x) = 4 \nextpassage
	&A(x)x \cdot \ddot{v}(x) + 2(\dot{A}(x)x + A(x)) \dot{v}(x) + (x \ddot{A}(x) + 2 \dot{A}(x) + 4xA(x)) v(x) = 4 \nextpassage
	&\dot{A}(x)x + A(x)=0 \to \ln(A(x)) = -\ln(x) \to A(x)= x^{-1}\nextpassage
	&\frac{1}{\cancel{x}} \cancel{x}\cdot \ddot{v}(x) + \left( \cancel{\cancel{x} \frac{2}{x^{\cancel{3}^2}}} - \cancel{\frac{2}{x^2}} + 4\cancel{x} \frac{1}{\cancel{x}} \right)v(x) = 4 \nextpassage
	&\triple{\ddot{v}(x) + 4v(x)= 4}{v \left( \frac{\pi}{4}\right)=0}{v \left( \frac{\pi}{2}\right)=0} \label{kekwdiff2}
\end{align}
Andiamo ora a risolvere l'\textbf{omogenea associata} dell'equazione in forma canonica
\begin{align}
	&\ddot{v}(x) + 4v(x) = 0 \firstpassage
	&\text{Ansatz } v(x)= e^{\alpha x} \nextpassage
	&\alpha^2 e^{\alpha} + 4 e^{\alpha} = 0 \nextpassage
	&\alpha^2 = -4 \to \alpha_{1,2} = \pm 2i \nextpassage
	&v_{{om}_1}(x) = e^{21x} \spacer v_{{om}_2}(x) = e^{-2ix} \nextpassage
	&v_{{om}}(x) = Ae^{2ix} + Be^{-2ix}
\end{align}
Dalla \ref{kekwdiff2} ricaviamo il nostro \textbf{operatore differenziale}
\begin{align}
	\mathcal{L}_x^c = \frac{d^2}{dx^2} + 4
\end{align}
Andiamo ora a verificare se è invertibile o meno, in base alle nostre Boundary Conditions:
\begin{align}
	&\double{v_{om}\left( \frac{\pi}{4}\right) =& Ae^{i\frac{\pi}{2}} + Be^{-i\frac{\pi}{2}}=iA-iB = 0}{v_{om}\left( \frac{\pi}{2}\right) =& Ae^{i\pi} + Be^{-i\pi} =-A -B = 0} \to \double{A-B=0}{A+B=0} \to \double{A=0}{B=0}\firstpassage
	&\mathcal{L}_x^c v=0 \leftrightarrow v=0 \quad \to \text{L'operatore è invertibile}
\end{align}
Andiamo ora a trovare due soluzioni dell'omogenea che soddisfino le BC
\begin{align}
	&v_A(x) \left( \frac{\pi}{4}\right) =0 \to iA - iB = 0 \to A=B \to v_A(x) = A(e^{2ix} + e^{-2ix}) \\
	&v_B(x) \left( \frac{\pi}{2}\right) =0 \to A+b = \to A=-B \to v_B(x) = A(e^{2ix} - e^{-2ix})
\end{align}
Siccome non abbiamo condizioni su A, la scegliamo comoda a seconda del caso, possiamo quindi scrivere
\begin{align}
	&v_A(x) = \frac{1}{2}(e^{2ix} + e^{-2ix}) = \cos(2x) \\
	&v_B(x) = \frac{1}{2i}(e^{2ix} - e^{-2ix}) = \sin(2x)\firstpassage
	& v_{om}(x) = \tilde{A} \cos(2x) + \tilde{B} \sin(2x) 
\end{align}

Andiamo ora a calcolare le \textbf{soluzioni fondamentali dell'equazione}
\begin{align}
	&W(v_A,v_B) = \begin{vmatrix}
		\cos(2x) && \sin(2x)\\
		-2\cos(2x) && 2 \cos(2x)
	\end{vmatrix} = \dots = 2\\
	&\tilde{G}(x,y) = \tilde{A} \cos(2x) + \tilde{B} \sin(2x) + \frac{\delta (x-y)}{2} \left[ \cos(2y) \sin(2x) - \cos(2x) \sin(2y) \right]\\
	&\mathcal{L}_x^c \tilde{G}(x,y) = \delta (x-y)
\end{align}

Per questo problema di SL, la \textbf{funzione di Green} sarà quindi data da
\begin{align}
	G(x,y) = \double{\frac{\cos(2x) \sin(2y)}{2} \quad x<y}{\frac{\cos(2y) \sin(2x)}{2} \quad x>y}
\end{align}
Da cui otteniamo
\begin{align}
	v(x) &= \int_{\frac{\pi}{4}}^{\frac{\pi}{2}} dy \; G(x,y) \cdot 4 = \nonumber \\
	&= 2\int_{\frac{\pi}{4}}^{x} dy \;\cos(2x) \sin(2y)+ 2\int_{x}^{\frac{\pi}{2}} dy \; \cos(2y) \sin(2x) = ... = \nonumber\\
	&=1 - \sin(2x) + \cos(2x) \firstpassage
	u(x) &= \frac{v(x)}{x} = \frac{1 - \sin(2x) + \cos(2x)}{x}
\end{align}

\bigskip

VERIFICA DA COPIARE QUANDO SONO MENO ABBOTTATO (pag 7 del pdf)

\newpage

\subsection{Esempio 4: Problema di Cauchy non lineare (Adriano Chialastri)}
Risolvere il seguente \textbf{problema di Cauchy}
\begin{align}
	\triple{x^2 \ddot{f}(x) + 3x \dot{f}(x) + f(x) = 0}{\dot{f}(1)=0}{f(1) =1}
\end{align}
Affinché l'operatore sia lineare dobbiamo avere BC omogenee, che non è il nostro caso. Dobbiamo quindi \textbf{linearizzare} il problema. Andiamo a definire
\begin{align}
	&u(x) = f(x) + mx + q\firstpassage
	&\dot{u}(x) = \dot{f}(x) + m
\end{align}	
Da cui otteniamo, imponendo che siamo omogenee le BC di $u(x)$
\begin{align}
	&\double{\dot{u}(1) =& \dot{f}(1) + m =0}{u(1) =& f(1) + m + q=0} \to \double{0 + m=0 }{1 + m + q=0} \to \double{m=& 0}{q=& -1} \firstpassage
	&u(x) = f(x) -1 \quad \leftrightarrow \quad f(x) = u(x)+1
\end{align}	
Il nostro problema diventa quindi
\begin{align}
	\triple{x^2 \ddot{u}(x) + 3x \dot{u}(x) + u(x) = -1}{\dot{u}(1)=0}{u(1) = 0}
\end{align}
Riscriviamo il problema in \textbf{forma canonica}
\begin{align}
	&u(x) = A(x)v(x)\firstpassage
	&\dot{u}(x) = \dot{A}(x) v(x) + A(x) \dot{v}(x)\nextpassage
	&\ddot{u}(x) = \ddot{A}(x) v(x) + 2\dot{A}(x) \dot{v}(x) + A(x) \ddot{v}(x)\nextpassage
	&[A(x) x^2] \ddot{v}(x) + [2\dot{A}(x) x^2 + 3A(x)x ]\dot{v}(x) + [\ddot{A}(x)x^2 + 3\dot{A}(x)x + A(x)]v(x) = -1 \nextpassage
	&2\dot{A}(x) x^2 + 3A(x)x=0 \to \dot{A}(x) = -\frac{3}{2}x^{-1}A(x) \to \ln(A(x)) = -\frac{3}{2} \ln(x)  \to A(x)= x^{-\frac{3}{2}} \nextpassage
	&\dot{A}(x) = -\frac{3}{2}x^{-\frac{5}{2}} \spacer \dot{A}(x) = -\frac{15}{2}x^{-\frac{7}{2}} \nextpassage
	&[x^{-\frac{3}{2}} x^2] \ddot{v}(x) + \left[x^2\left( \frac{15}{4}x^{-\frac{7}{2}} \right) + 3x \left( -\frac{3}{2} x^{-\frac{5}{2}} \right) + x^{-\frac{3}{2}} \right] v(x) = -1 \nextpassage
	&[x^{-\frac{3}{2}} x^2] \ddot{v}(x) + \left[\frac{15}{4} -\frac{9}{2} +1 \right] x^{-\frac{3}{2}} v(x) = -1 \nextpassage
	&x^{\frac{1}{2}} \ddot{v}(x) + \frac{1}{4}x^{-\frac{3}{2}} v(x) = -1 \nextpassage
	&\triple{\ddot{v}(x) + \frac{1}{4x^2} \dot{v}(x)  = -x ^{-\frac{1}{2}}}{\dot{v}(1)=0}{v(1) = 0}
\end{align}
Andiamo ora a risolvere l'\textbf{omogenea associata} dell'equazione in forma canonica
\begin{align}
	\ddot{v}(x) + \frac{1}{4x^2} \dot{v}(x)  = 0
\end{align}
Siccome è un'equazione di Eulero usiamo l'Ansatz $v(x) = x^\alpha$ e otteniamo
\begin{align}
	&\alpha(\alpha-1) + \frac{1}{4} = 0 \to \alpha_{1,2} = \frac{1\pm \sqrt{1-1}}{2} = \frac{1}{2} \firstpassage
	&v_{{om}_1}(x) = x^{\frac{1}{2}} \spacer v_{{om}_2}(x) = x^{\frac{1}{2}} \ln(x) \nextpassage
	&v_{om}(x) = Ax^{\frac{1}{2}} + Bx^{\frac{1}{2}}\ln(x) = x^{\frac{1}{2}}(A + B\ln(x)) \firstpassage
	&W(v_1,v_2) = \begin{vmatrix}
		x^{\frac{1}{2}} && x^{\frac{1}{2}}\ln(x)\\
		\frac{1}{2} x^{-\frac{1}{2}} && \frac{1}{2} x^{-\frac{1}{2}} \ln(x) + x^{-\frac{1}{2}}
	\end{vmatrix} = \dots = 1
\end{align}
La soluzione generica del PdC sarà data da
\begin{align}
	\tilde{G}(x,y) = \alpha v_1(x) + \beta v_2(x) + \theta(x-y)\left[ \frac{v_1(y) v_2(x) -v_1(x)v_2(y)}{W} \right] 
\end{align}
Siccome siamo in presenza di un PdC, a noi serve la \textbf{funzione di Green avanzata}, e quindi dobbiamo porre $\alpha=\beta=0$
\begin{align}
	G(x,y) &= \theta(x-y)\left[ \frac{v_1(y) v_2(x) -v_1(x)v_2(y)}{W} \right] = \nonumber \\
		   &= \theta(x-y)\left[ y^{\frac{1}{2}} x^{\frac{1}{2}} \ln(x) -x^{\frac{1}{2}}y^{\frac{1}{2}} \ln(y) \right] 
\end{align}
Possiamo quindi andare a calcolare la soluzione
\begin{align}
	v(x) &= \int_{1}^{x} dy \; G(x,y) \left( -y^{-\frac{1}{2}} \right) + \cancel{BC} = \nonumber\\
		 &= -\int_{1}^{x} dy \; \delta(x-y)\left[ x^{\frac{1}{2}} \ln(x) -x^{\frac{1}{2}} \ln(y) \right]= \nonumber\\
		 &= -x^{\frac{1}{2}}\int_{1}^{x} dy \; \left[ \ln(x) - \ln(y) \right] = \nonumber\\
		 &= -x^{\frac{1}{2}} \left[\ln(x)\int_{1}^{x} dy + \int_{1}^{x} dy\; \ln(y)\right] = \nonumber\\
		 &= -x^{\frac{1}{2}} \left[\ln(x)(x-1) - [y(\ln(y) -1)]_1^x \right] = \nonumber\\
		 &= -x^{\frac{1}{2}} \left[\ln(x)(x-1) - x(\ln(x) -1) -1 \right] = \nonumber\\
		 &= -x^{\frac{1}{2}} \left[\cancel{x\ln(x)} -\ln(x) - \cancel{x\ln(x)} +x -1 \right] \nonumber\\
		 &= x^{\frac{1}{2}} \left[1 +\ln(x)-x \right] \firstpassage
	u(x) &= x^{-\frac{3}{2}} v(x) = \frac{1}{x} \left[1 +\ln(x)-x \right] \firstpassage
	f(x) &= u(x) + 1 = \frac{1 + \ln(x)}{x}
\end{align}
Verifichiamo la bontà della soluzione trovata
\begin{align}
	&\dot{f}(x) = -\frac{\ln(x)}{x^2} \spacer \ddot{f}(x) = \frac{2\ln(x) -1}{x^3} \firstpassage
	&\double{\dot{f}(1) = -\frac{\ln(1)}{1}=0 }{f(1) = \frac{1 + \ln(1)}{1} = 1} \quad \leftarrow \quad \text{le BC sono rispettate}\\
	&\cancel{x^2}\frac{2\ln(x) -1}{x^{\cancel{3}}} - 3\cancel{x} \frac{\ln(x)}{x^{\cancel{2}}} + \frac{1+\ln(x)}{x} = 0 \to \frac{1}{x}[\cancel{2\ln(x)} -\cancel{3\ln(x)} + \cancel{\ln(x)} -\cancel{1} +\cancel{1}] = 0 \quad \leftarrow \quad \text{ok}
\end{align}
\newpage 
\subsection{Esempio 5: prova d'esame (Alessandro Marcelli)}
Sia il seguente problema di Cauchy
\begin{align}
	\triple{x^2 \ddot{f}(x) - x \dot{f}(x) + f(x) = x^2}{\dot{f}(1) = 0}{f(1)=0}
\end{align}
Si richede di calcolare
\begin{enumerate}
	\item la funzione di Green del problema
	\item la soluzione della non omogenea col metodo di Green per $x\geq 1$
\end{enumerate}
Iniziamo passando l'equazione in \textbf{forma omogenea} (omettiamo la dipendenza da x per snellire la notazione)
\begin{align}
	&f(x) = A(x)v(x) \to f=Av\firstpassage
	&\dot{f} = \dot{A}v + A \dot{v} \to \ddot{f} = \ddot{A}v + 2 \dot{A}\dot{v} + A \ddot{v} \nextpassage
	&x^2 \ddot{A}v + 2x^2 \dot{A}\dot{v} +x^2 A \ddot{v} -x\dot{A}v -x A \dot{v} + Av = x^2\nextpassage
	& x^2 \ddot{A}v +x(2x \dot{A} - A) \dot{v} + (x^2 \ddot{A} - x \dot{A} + A)v = x^2\nextpassage
	&2x \dot{A} - A = 0 \to \frac{dA}{dx} = \frac{A}{2x} \to \int dA \; \frac{1}{A} = \int dx \; \frac{1}{2x} = \ln(A) = \ln(\sqrt{x}) \nextpassage
	&A = \sqrt{x} = x^{\frac{1}{2}} \to \dot{A} = \frac{1}{2}x^{-\frac{1}{2}} \to \ddot{A} = -\frac{1}{4}x^{-\frac{3}{2}} \nextpassage
	&x^{\frac{5}{2}} \ddot{v} + \frac{1}{4} x^{\frac{1}{2}} = x^2\nextpassage
	&\ddot{v} +\frac{1}{4} x^{-2} v = x^{-\frac{1}{2}}
\end{align}
Cominciamo col risolvere l'omogenea:
\begin{align}
	\ddot{v} +\frac{1}{4} x^{-2} v = 0 
\end{align}
Siamo in presenza di un'equazione di eulero, visto che
\begin{align}
	&\ddot{v}  + \frac{a}{x} \dot{v} + \frac{b}{x^2} u = 0 \spacer a=0 \; , \; b= \frac{1}{4} \firstpassage
	&	\ddot{v} +\frac{1}{4x^2} v = 0 
\end{align}
Possiamo quindi fare l'Ansatz $v = x^\alpha$, ottenendo
\begin{align}
	\alpha (\alpha -1) x^{\alpha -2} +\frac{1}{4} x^{\alpha -2} = 0 \to \alpha^2 - \alpha + \frac{1}{4} = 0 
\end{align}
Siccome $\Delta = 0$ avremo una sola soluzione in $\alpha = \frac{1}{2}$. Ci conviene quindi trasformare l'equazione con il cambio di variabile $x=e^{t}$, dove otteniamo
\begin{align}
	&\frac{dv}{dx} = e^{-t}\frac{dv}{dt} \to \frac{d^2v}{dx^2} = e^{-2t}\frac{d^2 v}{dt^2} - e^{-2t}\frac{dv}{dt} \firstpassage
	&e^{-2t} \left[ \frac{d^2 v}{dt^2} - (a-1) \frac{dv}{dt} + bv\right] = 0
\end{align}
nel nostro caso $a=0$, quindi
\begin{align}
	\frac{d^2 v}{dt^2} + \frac{dv}{dt} + \frac{v}{4}= 0
\end{align}
Il nostro  Ansatz diventa $v = x^\alpha = (e^t)^\alpha  = e^{\alpha t}$ eotteniamo le due soluzioni
\begin{align}
	v_1 = e^{\alpha t} \spacer v_2 = t e^{\alpha t} 
\end{align}
Che, tornando nella variabile originale diventano
\begin{align}
	v_1 = x^\alpha \spacer v_2 = x^\alpha \ln(x) 
\end{align}
Nel nostro caso saranno quindi
\begin{align}
	&v_1 = x^\frac{1}{2} \spacer v_2 = x^\frac{1}{2} \ln(x) \firstpassage
	&v_{om} = x^\frac{1}{2}(A + B\ln(x)) \firstpassage
	&\dot{v}_{om} = \frac{1}{2}x^{-\frac{1}{2}}A + B \left( x^\frac{1}{2} \cdot \frac{1}{x} + \frac{1}{2}x^{-\frac{1}{2}} \right) = \frac{1}{2}x^{-\frac{1}{2}} A + \frac{3}{2}x^{-\frac{1}{2}} B = \frac{1}{2}x^{-\frac{1}{2}} (A + 3B) \nextpassage
	&\double{v_{om}(1) = 0}{\dot{v}_{om}(1) = 0} \to \double{1^\frac{1}{2}(A + B\ln(1)) = 0 \to& A = 0}{\frac{1}{2}\cdot 1^{-\frac{1}{2}} (A + 3B)=0 \to& B = 0}
\end{align}

La soluzione omogenea è nulla, e posso quindi usare il \textbf{teorema di Green}.
 
Iniziamo calcolando il \textbf{Wronskiano}
\begin{align}
	 W = \begin{vmatrix}
	x^\frac{1}{2} && x^\frac{1}{2} \ln(x) \\
	\frac{1}{2}x^{-\frac{1}{2}} && \frac{1}{2}x^{-\frac{1}{2}} \ln(x) + x^{-\frac{1}{2}}
\end{vmatrix} = \frac{1}{2} \ln(x) + 1 - \frac{1}{2}\ln(x) = 1
\end{align}
Andiamo ora calcolare la \textbf{funzione di Green avanzata} del problema
\begin{align}
	G(x,y) &= \theta(x-y)\left[ \frac{v_1(y) v_2(x) -v_1(x)v_2(y)}{W} \right] = \nonumber \\
	&= \theta(x-y)\left[y^\frac{1}{2} x^\frac{1}{2} \ln(x) - x^\frac{1}{2} y^\frac{1}{2}\ln(y) \right] 
\end{align}
Possiamo ora andare a calcolare la soluzione particolare
\begin{align}
	v(x) &= \int_{1}^{x} dy \; G(x,y) \cdot y^2 = \nonumber\\
		 &= \int_{1}^{x} dy \; \left[y^\frac{1}{2} x^\frac{1}{2} \ln(x) - x^\frac{1}{2} y^\frac{1}{2}\ln(y) \right] y^2 = \nonumber\\
		 &= x^\frac{1}{2} \ln(x)\int_{1}^{x} dy \;y^\frac{5}{2} - x^\frac{1}{2}\int_{1}^{x} dy \; y^\frac{5}{2}\ln(y)  = \nonumber\\
		 &= x^\frac{1}{2} \ln(x) \left.\frac{2}{7}y^{\frac{7}{2}} \right|_1^x + 
\end{align}

\newpage
\section{Equazioni Differenziali Complesse}

N'aggio fatto in tempo, tanto Bufalini le ha scritte bene su Latex.


\backmatter
% bibliography, glossary and index would go here.

\end{document}